\documentclass[11pt,a4paper]{article}
\usepackage[T1]{fontenc}
\usepackage[utf8]{inputenc}
\usepackage{lmodern}
\usepackage[margin=27mm]{geometry}
\usepackage{amsmath,amssymb,amsthm,mathrsfs}
\usepackage{booktabs,longtable,array}
\usepackage{graphicx}
\usepackage{hyperref}
\usepackage{xurl}
\usepackage{microtype}
\hypersetup{colorlinks=true,linkcolor=blue,urlcolor=blue}
\setlength{\emergencystretch}{3em}
\providecommand{\tightlist}{\setlength{\itemsep}{0pt}\setlength{\parskip}{0pt}}
\providecommand{\pandocbounded}[1]{#1}
\newtheorem{theorem}{Theorem}[section]
\newtheorem{proposition}[theorem]{Proposition}
\newtheorem{lemma}[theorem]{Lemma}
\newtheorem{corollary}[theorem]{Corollary}
\theoremstyle{definition}
\newtheorem{definition}[theorem]{Definition}
\newtheorem{remark}[theorem]{Remark}
\title{DEG\_L: Degrees, theta norms, and semiregularity\\
in the 032 construction}
\author{DEG\_L research run\\GPT-6 Astra Pro}
\date{21 November 2026}
\begin{document}
\maketitle
\begin{abstract}
This working paper separates the effective input surface, its nonzero mixed
period, and the signed codimension-two output of the 032 construction.
It proves exact Chow cancellations, integral norm-and-trace division,
fixed-field Weil projectors, and quantitative extraction of any fixed Weil
class from an effective witness. The theta spanning theorem supplies an
exact finite character-theta formula for the energy, and an explicit
surface-degree inequality. The exact quartic spectral formula is being
developed with its full source representation and measure data retained.
The original construction does not specify a smooth embedding; universal
semiregularity of its deeper-level morphism presentations is false. The
remaining deformation problem is formulated for a precise replacement
object and its full obstruction trace. A proper Hilbert parameter-space
argument carries a uniform effective witness bound on dense fibres to
bounded division on the entire family. The general uniform witness bound
in the open Weil families is not asserted. This is a live checkpoint
manuscript; the session has not ended.
\end{abstract}
\tableofcontents
\clearpage
\section{Scope, assumptions, and check protocol}
The prompt and the attached bridge sketch determine the questions. The
sketch is treated as a map of questions, not as a source of proved facts.
As requested, the attached 032 CM Hodge theorem and the stated Alp\"oge
Dirichlet GRH theorem are admitted inputs. The full Lean dependency archive
was not supplied or compiled. General Hecke GRH is a conditional input
attributed by the prompt to Masullo, \S1.2; the attached index does not
contain a complete numbered Hecke-GRH statement. No trace-formula number
is silently substituted for it.

Every finding has a stable run identifier: T6.x, Bx, S-xx, or DIVx within
DEG\_L\_20261121. Findings are formulations proved in this run, without a
claim that all elementary ingredients are new to the literature. A
conditional implication is proved only with its displayed hypotheses.
The detailed event ledger distinguishes whole-proof reading, targeted
checks, background use, and computations. Exact reasoning effort is not
exposed by the execution environment and is not invented.

The cycle equality $Nt=[Z^+]-[Z^-]$ means equality in cohomology under the
cycle-class map, unless a Chow-group equality is explicitly written. The
word degree always includes the chosen polarization and distinguishes a
pushforward cycle with multiplicity from its reduced image.

\part{Checked findings and explicitly conditional reductions}

\section{The surface degree and the Rallis comparison}

Date: 2026-11-21. Actor: GPT-6 Astra Pro, effort not exposed, child run
\texttt{/root/rallis\_t6}. This report records checked reductions, not a
claimed proof of T6 as originally phrased. The later T6.9 strengthens
the early analytic gate: all needed one-form theta expansions do exist
by the supplied human proof and checked literature; their coefficients
and integral normalization remain to be bounded.

\subsection{1. What the attached 032 construction actually
supplies}\label{what-the-attached-032-construction-actually-supplies}

Source:
\texttt{OPENAI\_032\_The-rational-Hodge-conjecture-for-CM-abelian-varieties.pdf},
Proposition 3.1, printed pp.~11--12, equations (3.1)--(3.5); Corollary
3.2, p.~12; Proposition 4.6, p.~23; Corollary 4.7, pp.~24--25, equations
(4.18)--(4.19); Corollary 6.7, p.~45; proof of Theorem 2.5, p.~46,
equation (7.1).

Write \(A=\prod_{i=1}^4B_i\), with each \(B_i=B(v_i)\) of dimension
\(s\), and \(n=4s\). In the nontrivial surface construction \(s\ge2\),
so \(n\ge8\). The surface \(S\) is a smooth projective ball quotient.
The paper obtains rational Hodge maps
\(h_i:H^1(B_i,\mathbf Q)\to H^1(S,\mathbf Q)\) and integers \(m_i>0\) so
that \(f_i:S\to B_i\) has \(f_i^*=m_i h_i\). It defines {[}
f=(f\_i),\quad z=f\_*{[}S{]}\in CH\^{}\{n-2\}(A),\quad k=\int\_S
f\_1\^{}*\alpha\_1,f\_2\^{}*\alpha\_2,f\_3\^{}*\alpha\emph{3,f\_4\^{}\emph{\alpha\emph{4\ne0.
{]} The generators \(\alpha_1,\alpha_2\) are holomorphic and
\(\alpha_3,\alpha_4\) antiholomorphic; their Betti coefficients are
algebraic. With copied factors \(B_i'\) and divisor kernels \(D_i\)
having pure class \(\alpha_i\otimes\beta_i\), the correspondence
\(K=\prod_i q_i^*D_i\) gives {[} \operatorname{cl}(K}}z)=k,t}\beta,
\qquad t\_\beta=\beta\_1\otimes\beta\_2\otimes\beta\_3\otimes\beta\_4.
{]} Thus three different quantities must be distinguished: the degree of
the effective pushforward surface \(z\); a mixed period \(k\) detecting
a switch line; and the codimension-two output \(K_*z/k\) with algebraic
coefficients. They have different degrees, codimensions, and positivity
properties. The paper itself explicitly says that \(z\) is a functional
on the line, not a vector on it. It only proves \(\dim f(S)=2\), not
that \(f\) embeds \(S\).

The theta construction first supplies nonzero total two-forms and their
mixed pairing. Section 7 then expands finite wedges, chooses a nonzero
summand, expands each theta class in rational CM-map images, and chooses
another nonzero summand. The final \(k\) is not stated to be the total
theta norm. No Rallis formula is stated in the attached 032 paper.

\subsection{2. Checked findings}\label{checked-findings}

\subsubsection{DEG-L.T6.1 --- Exact geometric degree and its positive
Gram
expression}\label{deg-l.t6.1-exact-geometric-degree-and-its-positive-gram-expression}

Let \(H\) be an ample line bundle on a complex abelian variety \(A\) of
dimension \(n\), and let \(f:S\to A\) be a morphism from a smooth
projective surface with two-dimensional image. Put \(z=f_*[S]\). Choose
translation-invariant holomorphic one-forms \(\omega_1,\ldots,\omega_n\)
such that the invariant positive representative of \(c_1(H)\) is {[}
\Omega\emph{H=\{i\over2\}\sum}\{a=1\}\^{}n\omega\_a\wedge\overline{\omega_a}.
{]} For \(\eta_a=f^*\omega_a\) and \(\xi_{ab}=\eta_a\wedge\eta_b\), one
has {[} \deg\_H(z)=\int\_S(f\^{}*\Omega\emph{H)\^{}2
=\{1\over2\}\sum}\{a\textless b\}\int\emph{S\xi}\{ab\}\wedge\overline{\xi_{ab}}.
\tag{G1} {]} If \(e=[\mathbf C(S):\mathbf C(f(S))]\), then \(z=e[f(S)]\)
and \(\deg_H(f(S))=\deg_H(z)/e\).

\textbf{Proof.} The projection formula gives the first equality. A
positive definite Hermitian form can be diagonalized, giving the
displayed expression for \(\Omega_H\). Expanding its square and moving
each antiholomorphic factor past the next holomorphic factor yields {[}
(f\^{}*\Omega\emph{H)\^{}2=\{1\over4\}\sum}\{a,b\}\eta\_a\wedge\eta\_b\wedge\overline{\eta_a}\wedge\overline{\eta_b}.
{]} The terms with \(a=b\) vanish; the terms with \((a,b)\) and
\((b,a)\) agree. This proves (G1). In local holomorphic coordinates, if
\(\xi=u\,dz_1\wedge dz_2\), then
\(\xi\wedge\bar\xi=4|u|^2dx_1dy_1dx_2dy_2\), so every summand is
nonnegative. The assertion about \(e\) is the definition of proper
pushforward of an integral cycle. QED.

\textbf{Check.} Read the entire proof, independently recomputed both
exterior-algebra signs and the factor \(1/2\), and compared \(A=S\) of
dimension two with the standard Euclidean Kähler form. In that test,
\(\Omega_H^2=2\,dV\) and \(\xi_{12}\bar\xi_{12}=4\,dV\).

\textbf{Meaning.} A degree is a positive sum of norms of all
polarization-coordinate wedges. It is not the selected mixed period of
four CM lines. Even when Rallis computes certain wedge norms, it takes
additional data to identify every wedge in (G1), its automorphic
spectral projections, and the conversion from invariant forms to
algebraic Betti generators.

\subsubsection{DEG-L.T6.2 --- Zero degree of the pure switch
class}\label{deg-l.t6.2-zero-degree-of-the-pure-switch-class}

Give \(A=\prod_{i=1}^4B_i\) the product polarization
\(H=\bigotimes_i\operatorname{pr}_i^*H_i\), where each \(H_i\) is ample.
Every class {[}
w\in H\textsuperscript{1(B\_1)\otimes H}1(B\_2)\otimes H\textsuperscript{1(B\_3)\otimes H}1(B\_4)\subset H\^{}4(A)
{]} has {[} \int\_A w,c\_1(H)\^{}\{n-2\}=0. \tag{G2} {]} Consequently a
nonzero rational switch class cannot be the cohomology class of an
effective codimension-two cycle for this polarization. A nonzero pure
tensor \(t_\beta\) also satisfies \(t_\beta^2=0\).

\textbf{Proof.} Expand \(c_1(H)^{n-2}\) into products
\(\prod_i\operatorname{pr}_i^*c_1(H_i)^{a_i}\). In factor \(i\), the
product with \(w\) has degree \(1+2a_i\), which is odd. Integration on
\(B_i\), of even real dimension \(2\dim B_i\), therefore vanishes. Every
term vanishes. Every nonzero effective cycle has strictly positive
degree against an ample polarization, so (G2) rules out effectivity. For
the last assertion, squaring a pure tensor repeats the same degree-one
class in every factor, and the square of a degree-one class vanishes in
characteristic zero. QED.

\textbf{Check.} Entire proof reread; parity argument checked separately
for unequal factor dimensions and arbitrary sums of pure tensors. It
does not say that a rational combination of conjugate switch lines has
zero polarized self-pairing; that would be false. The statement about
\(t_\beta^2\) concerns one pure algebraic-coefficient tensor only.

\textbf{Meaning.} The output of 032, Proposition 3.1, has degree zero
for a natural ample polarization. Therefore an equality between its
degree and a nonzero positive theta norm cannot hold. The complexity
\(D\) in Route 2 is the separate degree of positive and negative
effective representatives; (G2) does not bound it. For the input surface
\(z\), the ordinary square is already zero by codimension when
\(n\ge8\), since \(2(n-2)>n\). A numerical self-intersection would
require a different pairing and an explicit dualizing correspondence.

\subsubsection{DEG-L.T6.3 --- Level alone does not determine the raw
surface
degree}\label{deg-l.t6.3-level-alone-does-not-determine-the-raw-surface-degree}

Fix every datum through the rational maps \(h_i\) and the source surface
\(S\) in 032, Proposition 3.1. Replacing all clearing integers \(m_i\)
by \(Mm_i\) and \(f_i\) by \([M]\circ f_i\) is still an allowed
construction on exactly the same source level and with the same initial
theta inputs. If \(H\) is symmetric, then {[}
\deg\_H(z\_M)=M\^{}4\deg\_H(z),\qquad k\_M=M\^{}4k. \tag{G3} {]} Thus
there is no degree of the raw pushforward surface determined by the
level alone under the choices specified in 032.

\textbf{Proof.} Multiplication by \(M\) acts as \(M\) on \(H^1(A)\) and
as \(M^2\) on \(H^2(A)\), so \([M]^*c_1(H)=M^2c_1(H)\). Applying the
projection formula twice gives
\(\deg_H([M]_*z)=\int_S f^*[M]^*c_1(H)^2=M^4\deg_H(z)\). Each of the
four factors in the period acquires \(M\). Positivity in T6.1 and the
nonzero period guarantee \(\deg_H(z)>0\). QED.

\textbf{Check.} Entire proof reread. The assertion concerns \(f_*[S]\),
the cycle defined by 032, not the reduced image \([f(S)]\); the generic
mapping degree can change. Even a symmetric assumption is unnecessary
for the cohomological degree identity, but retained for a standard
polarization normalization.

\subsubsection{DEG-L.T6.4 --- Auxiliary-cover division cancels in
Chow}\label{deg-l.t6.4-auxiliary-cover-division-cancels-in-chow}

Keep the target \(A\), kernels \(K\), and algebraic generators fixed.
For a finite surjective morphism \(\pi:S'\to S\) of degree \(d\) between
smooth projective surfaces, use \(f'=f\pi\) and the pulled one-forms.
Then {[} z'=d z,\quad k'=d k,\quad {K_*z'\over k'}=\{K\_*z\over k\}
\quad\hbox{in }CH\^{}2(A')\otimes\overline{\mathbf Q}. \tag{G4} {]} In
particular this applies to every deeper principal cover used in 032.
Every fixed rational coefficient-descent combination gives one identical
rational Chow class throughout the tower. A single rational
representative therefore gives bounded \(N,D\) throughout this auxiliary
tower.

\textbf{Proof.} Proper pushforward gives \(\pi_*[S']=d[S]\), and
functoriality gives \(z'=f_*\pi_*[S']=dz\). The integration formula for
a finite map gives \(k'=dk\); alternatively use its projection formula
with the degree-four class on \(S\). The correspondence action is
linear, so the quotient is identical before taking any cycle-class map.
Coefficient trace and rational linear combination commute with this
identity. Once such a fixed rational cycle is written as a difference of
effective cycles after clearing denominators, the same \(N,D\) work at
every cover. QED.

\textbf{Check.} Entire proof reread; arbitrary finite surjective maps
work, although the 032 covers are finite étale. This proves division for
auxiliary covers at a fixed abelian target. It does not transport a
cycle to new targets in a dense Hecke orbit.

\subsubsection{DEG-L.T6.5 --- Clearing-map division also cancels with
biextension
kernels}\label{deg-l.t6.5-clearing-map-division-also-cancels-with-biextension-kernels}

The divisor kernels in Proposition 3.1 may be chosen as algebraic linear
combinations of birigidified Poincaré-type line bundles, without
changing their Betti classes. For these choices, write
\(u=\prod_i[M_i]\) on \(A\), with \(M_i>0\). Then {[}
(u\times1)\^{}*K=\Bigl(\prod\_i M\_i\Bigr)K
\quad\hbox{in }CH\^{}4(A\times A')\otimes\overline{\mathbf Q}. {]}
Replacing \(f_i\) by \([M_i]f_i\) leaves \(K_*z/k\) unchanged in Chow.

\textbf{Proof.} For a divisor class \(D\) on \(B\times B'\), replace any
representative line bundle \(L\) by its birigidified part {[}
L\^{}\{\mathrm b\}=L\otimes\operatorname{pr}\_1\^{}*(L\textbar\_\{B\times0\})\^{}\{-1\}
\otimes\operatorname{pr}\_2\^{}*(L\textbar\_\{0\times B'\})\^{}\{-1\},
{]} with the evident normalization of the fiber at \((0,0)\). Extend
this operation linearly to the algebraic-coefficient divisor kernel.
Since its prescribed Betti class is purely in \(H^1(B)\otimes H^1(B')\),
this operation leaves that class unchanged. A birigidified line bundle
corresponds to a morphism \(B'\to\widehat B\) sending zero to zero,
hence to a homomorphism, and is the pullback of the Poincaré bundle. Its
biextension property gives \(([M]\times1)^*D=M D\) in \(CH^1\).
Multiplying the four kernel identities yields the displayed formula.

The projection formula for \(u\times1\) then gives {[}
K\_\emph{u\_}z=p'\_\emph{\bigl(p\^{}}z\cdot(u\times1)\^{}\emph{K\bigr)
=\Bigl(\prod\emph{iM\_i\Bigr)K}}z. {]} The period is multiplied by
exactly the same product, because the four pullbacks on \(H^1\) multiply
by \(M_i\). Division cancels in Chow. QED.

\textbf{Check.} Entire proof reread. The factor is \(\prod_iM_i\), not
\(\prod_iM_i^2\). For common \(M\) this is \(M^4\), consistent with
T6.3. The permitted change of representatives may alter the original
unnormalized Chow class, but preserves its required Betti class and
gives a valid 032 construction. The claim is not that an arbitrary
representative of a pure cohomology divisor automatically has the exact
biextension scaling in Chow.

\textbf{Grothendieck locator.} SGA 7 I, Exposé VII, Example 2.9.5 and
Remark 2.9.6 (canonical biextension structure on the Poincaré bundle),
and VII (2.4.2) (additivity in the variables). Exact original scans need
to be read before final citation; primary PDF discovered at
\texttt{https://library.slmath.org/nonmsri/sga/sga/pdf/sga7-1.pdf}, but
web retrieval rejects its 37 MB size. Secondary cross-check of locators:
Pépin, \emph{Néron's pairing and relative algebraic equivalence},
Algebra \& Number Theory 6 (2012), §3.1, explicitly cites VII 2.9.5.
These are published SGA locators, not an invented manuscript cote/scan.

\subsubsection{DEG-L.T6.6 --- A mixed period gives a lower, not upper,
surface-degree
estimate}\label{deg-l.t6.6-a-mixed-period-gives-a-lower-not-upper-surface-degree-estimate}

Use a product polarization on \(A\). Let the four chosen forms in the
032 period be supported on their four distinct factors, and let \(r_i\)
be their norms in the invariant Hermitian cotangent metric for \(H\)
(use \(\bar\alpha_i\) for the last two holomorphic cotangent vectors).
Then {[} \textbar k\textbar{}\le (r\_1r\_2r\_3r\_4)\deg\_H(z). \tag{G5}
{]} There is no upper degree estimate from this single linear period
over the cone of effective surface cycles.

\textbf{Proof.} The four unit forms
\(\alpha_1/r_1,\alpha_2/r_2,\bar\alpha_3/r_3,\bar\alpha_4/r_4\) are
orthonormal because they lie on different polarized factors. Extend them
to an orthonormal cotangent basis in T6.1. Cauchy--Schwarz and
\(2ab\le a^2+b^2\) give {[} \{\textbar k\textbar{}\over\prod\emph{i
r\_i\}\le {1\over2}(\textbar{}\xi}\{12\}\textbar\^{}2+\textbar{}\xi\_\{34\}\textbar\^{}2)
\le\deg\_H(z). {]} For the final assertion, take a nonzero effective
complete-intersection surface \(Y\) with class a positive multiple of
\(c_1(H)^{n-2}\). Its pairing with the switch tensor is zero by T6.2,
while its degree is positive. The effective cycles \(z+mY\) all give the
same mixed period and have degrees tending to infinity. QED.

\textbf{Check.} Entire proof reread; this is a statement about all
effective cycles, not necessarily individual connected 032 ball-quotient
images. The common-source scaling counterexample T6.3 addresses the
actual construction's nonuniqueness separately. The lower bound depends
on metric norms/CM periods of the generators; an L-value estimate for
\(k\) alone does not eliminate them.

\subsection{3. What Rallis can actually compute, and the remaining
explicit
gate}\label{what-rallis-can-actually-compute-and-the-remaining-explicit-gate}

Primary source: W. T. Gan, Y. Qiu, S. Takeda, \emph{The regularized
Siegel--Weil formula (the second term identity) and the Rallis inner
product formula}, arXiv:1207.4709v3, §1.1, §11.6, Theorem 11.4, printed
pp.~54--55. URL: \texttt{https://arxiv.org/pdf/1207.4709}.

The theorem is about theta \textbf{integrals of specified automorphic
representations}, not a kernel evaluated at an unspecified source point.
For a unitary source of dimension \(n_0\) and target of dimension \(m\),
its normalization uses \(s_{m,n_0}=(m-n_0)/2\) and the standard value at
\(s_{m,n_0}+1/2\), retaining the splitting-character twist and
normalized doubling integrals. In the rank-two source / rank-three
target case it gives {[}
\langle\theta(\phi\_1,f\_1),\theta(\phi\_2,f\_2)\rangle
=2L(1,\pi\times\chi\_V),Z\^{}*(1/2,\phi\_1\otimes\bar\phi\_2,f\_1,f\_2),
\tag{R1} {]} when the theorem's hypotheses and conventions hold. Here
the target is globally anisotropic in 032, hence has Witt index zero, so
the lower-tower vanishing condition is empty. The scalar automorphic
inner product is converted exactly to the differential-form inner
product in T6.10 below, including component/level volume.

For the rank-one source / rank-three target, the formal point is \(3/2\)
in unitary normalization; it lies outside the displayed dimensional
range of Theorem 11.4. A convergent-range theorem must be used instead.
The bounds agent located, and this agent independently opened and
checked, J. Haan, \emph{The Bessel period of U(3) and U(2) involving a
non-tempered representation}, arXiv:1403.5061v7, §3.3, Theorem 3.4,
printed p.~8. In its conventions the character-lift norm is {[}
\langle\theta(\bar f\_1,\phi\emph{1),\theta(\bar f\_2,\phi\emph{2)\rangle
=\{L\_E(3/2,\mathrm{BC}(\tau)\otimes\gamma\^{}3)\over L\_F(3,\epsilon}\{E/F\})\}
\prod\emph{v Z\_v\^{}\#(1,f}\{1,v\},f}\{2,v\},\Phi\_\{1,v\}). {]} Its
§3.6.1, p.~9, states absolute convergence and nonvanishing of the
numerator Euler product at \(3/2\). Thus this particular global factor
does not require GRH. The version is identified in the PDF itself; use
the unversioned URL \texttt{https://arxiv.org/pdf/1403.5061}. Do not
cite GQT Theorem 11.4 for the rank-one/rank-three case.

The total theta tensor in 032, Proposition 4.6, uses a Hermitian plane
\(W\). Its dual partner is \(U(W)\), not \(U(1)\). Spectral projection
can therefore require a representation of a unitary group of rank two.
Its degree-two standard \(L\)-function is not automatically a
Hecke-character \(L\)-function. No claim is made here that a particular
nonabelian constituent must occur in the paper's chosen tensor: rather,
character-only decomposition is a specific additional theorem to be
proved, not a consequence of the mere existence of the theta kernel. All
local zeta factors, level volume, input scalars and CM-period
comparisons remain part of the equality.

The further attacks below complete the existential analytic comparison:
T6.9 supplies every one-form theta expansion, T6.10 fixes the
scalar/geometric normalization, T6.11 gives an exact finite
Hecke-\(L(3/2)\) energy formula and degree upper bound, and T6.12--T6.13
give an exact finite standard-\(L(1)\) degree expression with all
comparison/input data retained. What is still not supplied is an
effective bound for those actual integral realization/input data on a
varying target orbit. T6.7--T6.8 separately control period division and
signed correspondence degrees once the geometric and kernel data are
bounded.

The exact-cancellation findings T6.4 and T6.5 remove two avoidable
sources of apparent denominator/level growth \textbf{before} pursuing
that harder gate. They leave the meaningful quantities: primitive
integral CM maps, their polarizations and periods, the genuinely new
choices of source/correspondence, and their behavior when the abelian
target varies on a Hecke orbit.

\subsection{4. Approaches tried and their
disposition}\label{approaches-tried-and-their-disposition}

\begin{enumerate}
\def\labelenumi{\arabic{enumi}.}
\tightlist
\item
  Direct identification of the final 032 class with a positive Rallis
  norm: refuted for product polarization by T6.2.
\item
  Identifying the selected four-form period with geometric degree:
  refuted in that form; exact replacement is T6.1, and T6.6 shows the
  selected period is only a lower-bound witness.
\item
  Defining raw degree as a function only of principal level: refuted
  without further normalization by T6.3.
\item
  Treating source-cover degree growth as a genuine obstruction to
  division: removed by the exact Chow cancellation T6.4.
\item
  Treating enlarged denominator-clearing maps as a genuine obstruction:
  removed after a legitimate biextension choice by T6.5.
\item
  Applying a rank-one Rallis formula directly to a rank-two wedge
  tensor: not justified. T6.12--T6.13 instead prove a finite cuspidal
  rank-two spectral expansion and standard-\(L(1)\) identity. No
  character-only or single-value equality has been proved.
\item
  Bounding the surface degree using only one-form norms: positive
  alternative sent to the bounds agent. Subsequently strengthened by
  T6.9 below: existence of every needed theta expansion is established
  from the human proof, and T6.10 fixes the scalar/geometric measure
  conversion. Effective coefficient/integral and test-input bounds
  remain. Hodge index gives
  \(\deg_H(z)\le E_f^2/\operatorname{vol}(S)\), where
  \(E_f=\int_Sf^*\Omega_H\wedge\omega_S\). This could use rank-one theta
  decompositions, but the rational CM extraction does not by itself
  identify all polarization-coordinate pullbacks as theta lifts with
  controlled coefficients.
\end{enumerate}

\subsection{5. Source-check ledger}\label{source-check-ledger}

\begin{itemize}
\tightlist
\item
  032 Proposition 3.1 and Corollary 3.2, pp.11--12: read entire
  statements and proofs, equations and codimension checked in this run.
  Imported CM theorem itself is background, not independently reproved.
\item
  032 Proposition 4.6 and Corollary 4.7, pp.23--25: read entire proofs
  to check fixed input, finite decomposition, dependence on arbitrary
  constant-line trivializations, and covering formula. No independent
  proof of analytic theta construction claimed.
\item
  032 Theorem 5.1 proof and Corollary 6.7, pp.44--45; Theorem 2.5 proof,
  p.46: read to track extraction from complex theta classes into
  algebraic Betti vectors and then rational maps. Extraction supplies
  spans, not a normalized unique cycle or coefficient bound.
\item
  Gan--Qiu--Takeda §11.6 and Theorem 11.4: primary arXiv PDF read via
  web, with theorem range and normalized factors checked. This is
  use-and-recheck, not a reproof.
\item
  SGA 7 original locator: discovered but original scan not yet
  inspected; must be marked unchecked until root's scan check. No
  transcription correction has been established.
\end{itemize}

\subsection{6. Further checked attack: algebraic norm clears the period
without a lower
bound}\label{further-checked-attack-algebraic-norm-clears-the-period-without-a-lower-bound}

\subsubsection{DEG-L.T6.7 --- Integral coefficient normalization and
trace
division}\label{deg-l.t6.7-integral-coefficient-normalization-and-trace-division}

Retain the 032 hypothesis that \textbf{every scalar conjugate} of each
tensor \(\alpha_i\otimes\beta_i\) is of type \((1,1)\); one
algebraic-coefficient Hodge class without this hypothesis would not
suffice. After rescaling the chosen algebraic Betti generators, fix a
finite Galois field \(L/\mathbf Q\), of degree \(d\), containing all
their coefficients. They may be taken in
\(H^1(B_i,\mathbf Z)\otimes\mathcal O_L\) and
\(H^1(B_i',\mathbf Z)\otimes\mathcal O_L\). The divisor kernels can then
be chosen in \(CH^1(B_i\times B_i')\otimes\mathcal O_L\), their product
\(K\) in \(CH^4(A\times A')\otimes\mathcal O_L\), and the nonzero period
\(k\) lies in \(\mathcal O_L\).

Let \(N=|N_{L/\mathbf Q}(k)|\). For every \(a\in\mathcal O_L\) the cycle
{[}
C\_a=\operatorname{Tr}\emph{\{L/\mathbf Q\}\bigl(a(N/k)K}*z\bigr)\in CH\^{}2(A')
{]} is integral and satisfies {[}
\operatorname{cl}(C\_a)=N,\operatorname{Tr}\emph{\{L/\mathbf Q\}(a
t}\beta). \tag{G6} {]} Thus clearing the algebraic period needs
\textbf{no analytic lower bound} for \(|k|\).

More explicitly, write \(K=\sum_{j=1}^J a_j\Gamma_j\) with
\(a_j\in\mathcal O_L\) and integral Chow correspondences \(\Gamma_j\).
Put {[}
M=\max\emph{\sigma\textbar{}\sigma k\textbar,\qquad B\_j=\max}\sigma\textbar{}\sigma(a
a\_j)\textbar. {]} Then \(M\ge1\), \(N\le M^d\), and the integral
coefficient {[} c\_j=\operatorname{Tr}\_\{L/\mathbf Q\}\bigl(a a\_j
N/k\bigr) {]} satisfies \(|c_j|\le d B_jM^{d-1}\). If
\(\Gamma_{j*}z=W_j^+-W_j^-\) with effective cycles of degrees at most
\(D_j\), then \(C_a\) admits effective positive and negative parts of
degrees at most {[} dM\^{}\{d-1\}\sum\_j B\_jD\_j. \tag{G7} {]}

\textbf{Proof.} The finitely many algebraic coordinates of the
generators have a common integer denominator, so multiplying the
generators clears it. For each \(i\) the class
\(\alpha_i\otimes\beta_i\) is then an integral \(\mathcal O_L\)-linear
combination of Betti classes. Expand it in an integral basis of
\(\mathcal O_L\). Its coefficient classes in integral cohomology are all
of type \((1,1)\): scalar conjugation preserves the Hodge property of
the whole tensor, and the embedding matrix of the integral basis is
invertible, so each coefficient is a complex linear combination of these
conjugate \((1,1)\) classes. The integral Lefschetz \((1,1)\) theorem
supplies line bundles for the integral coefficient classes. Choosing
their birigidified parts as in T6.5 preserves the required class. This
produces the integral-coefficient kernels. Pullback by actual morphisms,
cup product and evaluation on the integral fundamental class show
\(k\in\mathcal O_L\).

The monic characteristic polynomial of the algebraic integer \(k\) shows
\(N/k\in\mathcal O_L\) (equivalently, up to sign it is the product of
all conjugates except the identity embedding). Hence every \(c_j\) is an
integer, and taking traces of \(\operatorname{cl}(K_*z)=kt_\beta\)
proves (G6). Since \(N\) is a nonzero integer and \(N\le M^d\), one has
\(M\ge1\). For each embedding \(\sigma\), {[}
\textbar{}\sigma(N/k)\textbar=\prod\_\{\tau\ne\sigma\}\textbar{}\tau(k)\textbar{}\le M\^{}\{d-1\}.
{]} Triangle inequality for the trace gives the coefficient estimate.
Replace \(c_jW_j\) by \(|c_j|W_j^+-|c_j|W_j^-\) if \(c_j\ge0\), and
interchange signs if \(c_j<0\). Adding all positive and negative parts
proves (G7). QED.

\textbf{Check.} Entire proof reread, including integral coefficient
descent, the distinction between an integral basis of \(\mathcal O_L\)
and a basis of eigenvectors, and the trace bound when \(k\) lies in a
proper subfield. Repeated conjugates cause no problem. No lower bound
for any individual \(\sigma k\) was used.

\textbf{Scope.} The targets \(\operatorname{Tr}(a t_\beta)\) are
integral Hodge classes. An arbitrary prescribed rational class in their
rational span may require a further rational coefficient \(a\) and a
clearing integer. Bounding that integral spanning index is a remaining
lattice problem; it is not suppressed by this theorem. Rescaling the
generators changes their metric sizes and the kernels, so it is not a
uniformity argument by itself.

\subsubsection{DEG-L.T6.8 --- Explicit signed-degree bound for the four
divisor
kernels}\label{deg-l.t6.8-explicit-signed-degree-bound-for-the-four-divisor-kernels}

Let \(A,A'\) have dimension \(n\), with polarizations \(H,H'\), and put
\(\mathscr L=p^*H\otimes p'^*H'\) on \(A\times A'\). Let
\(D_1,\ldots,D_4\) be integral divisor classes. Choose expressions
\(D_i=D_i^+-D_i^-\) with \(D_i^\pm\) represented by basepoint-free
effective divisors, and positive integers \(b_i\) such that both
\(D_i^\pm\) and \(b_i c_1(\mathscr L)-D_i^\pm\) are nef. Such choices
exist for every fixed divisor class after tensoring with a sufficiently
high power of an ample bundle. Let \(z\) be an effective surface cycle
on \(A\). Then {[} (D\_1D\_2D\_3D\_4)\_*z=W\textsuperscript{+-W}-, {]}
with effective codimension-two cycles on \(A'\) satisfying {[}
\deg\_\{H'\}W\^{}\pm\le48\Bigl(\prod\_i b\_i\Bigr),(H'\^{}n),\deg\_H z.
\tag{G8} {]}

\textbf{Proof.} Expand the product of four differences into sixteen
signed products. For each product choose the four divisors successively
in their basepoint-free linear systems so that they meet each component
of \(z\times A'\) properly. The resulting intersection cycle is
effective; proper pushforward is effective, with components whose
dimension drops contributing zero. Its degree is the intersection of
that product with \(p'^*c_1(H')^{n-2}\) on \(z\times A'\).

Replacing one nef divisor \(D_i^\pm\) by \(b_i c_1(\mathscr L)\)
increases or preserves this intersection, since the difference and the
other divisors are nef. Iterating bounds the degree by {[}
\Bigl(\prod\emph{i b\_i\Bigr)
\int}\{z\times A'\}(p\^{}*c\_1(H)+p'\textsuperscript{\emph{c\_1(H'))\textsuperscript{4p'}}c\_1(H')}\{n-2\}.
{]} Only the term with two factors from \(A\) and two from \(A'\)
contributes. Its binomial coefficient is \(\binom42=6\), and its
integral is \(6\deg_H(z)(H'^n)\). There are eight terms of each sign.
Summing them proves (G8). QED.

\textbf{Check.} Entire proof reread; dimensions are
\(\dim(z\times A')=n+2\), four divisor intersections leave dimension
\(n-2\), and the final degree uses exactly \(n-2\) target hyperplanes.
Coefficient \(48=8\binom42\) checked. Nef intersection monotonicity is
used only on an effective cycle with all remaining factors nef. The
assertion is a bound for a selected signed representative, not for every
rationally equivalent cycle.

Combining T6.7 and T6.8 gives a completely explicit denominator/degree
bound from the source surface degree, the finitely many integral
divisor-kernel choices, their coefficient embeddings, the
coefficient-field degree, and the Galois conjugate periods. By T6.6
applied to each conjugate (reordering the two holomorphic and two
antiholomorphic factors as necessary), {[} M\le G\deg\emph{H z,\qquad
G=\max}\sigma\prod\_\{i=1\}\^{}4\textbar{}\sigma\alpha\_i\textbar\_H.
{]} For fixed \(L\), generators, and kernel data, the displayed bounds
are polynomial of degree at most \(d\) in \(\deg_H z\). This is an
actual sufficient criterion for bounded division if those data are
uniformly bounded on the chosen target orbit. It does not prove their
uniform boundedness. It removes the separate purported need for a lower
bound on the extracted algebraic period; the remaining nonvanishing
input is exactly the already established \(k\ne0\).

\subsection{7. Further checked attack: the human proof supplies every
one-form theta
expansion}\label{further-checked-attack-the-human-proof-supplies-every-one-form-theta-expansion}

\subsubsection{DEG-L.T6.9 --- Finite-level theta spanning removes the
existence
gate}\label{deg-l.t6.9-finite-level-theta-spanning-removes-the-existence-gate}

For a sufficiently deep principal ball-quotient surface in the 032
construction, \textbf{every} holomorphic one-form, hence every pullback
of a polarization coordinate under its actual map to an abelian variety,
is a finite complex linear combination of genuine integrated theta lifts
from automorphic characters of one-dimensional unitary groups. The
finite combination can be made at the chosen common principal level.
This is an existence statement; it does not give bounds on its
coefficients or on its comparison with the integral Betti/Hom lattice.

\textbf{Sources.} The supplied van Hoften--Sempliner paper, Proposition
3.1.7 and Theorem 3.1.11, printed p.11; proof of Proposition 4.1.8,
printed p.15, especially its explicit all-\(H^{1,0}\) spanning
assertion. Primary literature checked: Bergeron--Millson--Mœglin,
\emph{The Hodge conjecture and arithmetic quotients of complex balls},
Acta Mathematica 216 (2016), Theorem 9.1, pp.73--75, with the range
stated in §9.1, p.72; \(3(a+b)+|a-b|<2m\) holds for \(m=3\),
\((a,b)=(0,1)\) because \(4<6\). The older arXiv:1306.1515v2 has
different numbering, Theorem 9.3, pp.57--59. Also Y. Liu,
arXiv:2102.11518v2, Proposition 4.13, printed p.29, and Theorem 4.18,
pp.32--33. The actual printed hypothesis in Proposition 4.13 is
\(n\ge3\), confirmed by rendering the page: web extraction incorrectly
displays \(n>3\).

\textbf{Proof of the finite-level form needed here.} First pass, if
necessary, to a deeper principal level. For an integral isometry of the
Hermitian lattice, its determinant is an algebraic unit all of whose
complex absolute values are one. Hence it is a root of unity. The group
of roots of unity of \(E\) is finite, so a sufficiently deep principal
congruence condition kills every nontrivial determinant. At that level
the principal arithmetic groups for \(U(V)\) and \(SU(V)\) coincide. The
ball quotient in question is therefore a connected component of the
compact unitary Shimura variety at that level. The extra cover is
harmless for the normalized Chow construction by T6.4.

Extend a holomorphic one-form on this connected component by zero on the
other components. The cited spanning theorem on the direct-limit
cohomology writes its class as a finite sum of integrated cohomological
theta classes; in degree one the source has dimension one, so the source
representations are automorphic characters. Take a common deeper level
on which all finitely many inputs are fixed, and average over the
original compact open level subgroup. The class to be expressed is
already invariant, and equivariance of the theta lift moves this
averaging to its finite Schwartz input. Each such average is a finite
sum, since a smooth finite-adelic input has an open stabilizer. Thus
every summand can be taken at the original level. Restrict to the chosen
component. Equality of classes gives equality of holomorphic forms
because \(H^0(S,\Omega^1_S)\to H^1(S,\mathbf C)\) is injective on a
compact Kähler variety. QED.

\textbf{Check.} Entire argument reread. The pro-level span is used only
for a finite collection of vectors; compact averaging is finite on each
smooth input. The determinant argument uses norm-one \textbf{units}, not
arbitrary norm-one elements of \(E^\times\), which need not be roots of
unity. No claim that arbitrary arithmetic finite covers are principal
was used. It suffices to deepen within the principal tower already
authorized by 032. Published BMM theorem and its whole proof were read;
the classification theorem itself is cited background, not independently
reproved.

\textbf{Concrete energy consequence.} Let \(\omega_S\) be an invariant
Kähler form and let \(\eta_j=f^*\omega_j=\sum_a c_{ja}\theta_a\) be one
of these finite expansions. Define the geometric Gram form {[}
G\_\{ab\}=\int\_S
\{i\over2\}\theta\_a\wedge\bar\theta\_b\wedge\omega\_S. {]} Then the
energy is exactly \(E_f=\sum_j\sum_{a,b}c_{ja}\bar c_{jb}G_{ab}\). Each
full-Shimura-variety diagonal norm can be expressed by the rank-one
Rallis formula (Haan, Theorem 3.4), after its explicitly chosen
scalar-cohomological and measure normalization. Restricting to one
component only lowers that nonnegative norm, so with those full norms
denoted by \(Q_a\) one has {[}
E\_f\le\sum\_j\left(\sum\emph{a\textbar c}\{ja\}\textbar{}\sqrt{Q_a}\right)\^{}2,
\qquad \deg\emph{H f}*{[}S{]}\le{1\over\int_S\omega_S^2}
\left\{\sum\_j\left(\sum\emph{a\textbar c}\{ja\}\textbar{}\sqrt{Q_a}\right)\textsuperscript{2\right\}}\{!2\}.
\tag{G9} {]} The last inequality is the Hodge-index bound proved in the
bounds report. This is an unconditional finite-character-\(L(3/2)\)
comparison after specifying the geometric realization coefficients and
all local/measure factors. It is not an effective degree bound in the
level alone until those quantities are controlled. The actual existence
of the expansions, formerly left as a gate in the earlier analysis, is
now resolved.

A further exact identity locates what prevents one-form norms from
giving the degree itself. Put \(\beta=f^*c_1(H)\),
\(V=\int_S\omega_S^2\), and \(\beta_0=\beta-(E_f/V)\omega_S\). The class
\(\beta_0\) is primitive of type \((1,1)\), so with its Hodge norm
\(\|\beta_0\|^2=-\int_S\beta_0^2\), {[} \deg\emph{H
f}*{[}S{]}=\{E\_f\^{}2\over V\}-\textbar{}\beta\_0\textbar\^{}2.
\tag{G10} {]} This follows by expanding the square and using
\(\beta_0\cdot\omega_S=0\). Thus the precise missing datum for an
equality is the primitive degree-two component of the pulled
polarization. The rank-two theta/see-saw approach is relevant to this
remaining fourth-moment term; the rank-one Rallis norms already suffice
for the upper comparison (G9).

\subsection{8. Exact scalar normalization and the finite-character
energy
formula}\label{exact-scalar-normalization-and-the-finite-character-energy-formula}

\subsubsection{DEG-L.T6.10 --- The scalar/differential-form norm
conversion is
explicit}\label{deg-l.t6.10-the-scalardifferential-form-norm-conversion-is-explicit}

Let \(G=U(V)\) be the compact-quotient unitary group used above, let
\(K_f\) be neat, and put {[}
X\_K=G(F)\backslash(D\times G(\mathbf A\_f)/K\_f). {]} The notation
includes every connected component. Choose the invariant metric on the
two-dimensional ball so that an orthonormal holomorphic coframe
satisfies {[}
\omega=\{i\over2\}(e\_1\wedge\bar e\_1+e\_2\wedge\bar e\_2),\qquad dV=\{\omega\^{}2\over2\}.
{]} Normalize Haar measure on \(K_\infty\) to have mass one. Write the
archimedean Haar factor in the chosen Tamagawa factorization as
\(dg_\infty=c_\infty\,dV\,dk\), and set {[}
a\_K=c\_\infty\operatorname{vol}\emph{\{dg\_f\}(K\_f). {]} For a
holomorphic one-form \(\eta\), let
\(F_\eta(g)\in E=(\mathfrak p^{1,0})^*\) be its moving-frame automorphic
realization and \(F_{\eta,e}(g)\) its coefficient along any fixed unit
vector \(e\in E\). Then {[}
\int}\{X\_K\}\{i\over2\}\eta\wedge\bar\eta\wedge\omega
=\{2\over a\_K\}\int\emph{\{{[}G{]}\}\textbar F}\{\eta,e\}(g)\textbar\^{}2,dg.
\tag{G11} {]} For a holomorphic two-form \(\xi\), let \(F_\xi\) denote
its scalar coefficient along the unit frame \(e_1\wedge e_2\). Then {[}
\int\emph{\{X\_K\}\xi\wedge\bar\xi
=\{4\over a\_K\}\int}\{{[}G{]}\}\textbar F\_\xi(g)\textbar\^{}2,dg.
\tag{G12} {]} Consequently, extend each \(\eta_j=f^*\omega_j\) by zero
from the chosen component \(S\) to \(X_K\). Its wedge extension is also
zero on the other components, and the degree formula becomes {[}
\deg\emph{H
f}*{[}S{]}=\{2\over a\_K\}\sum\emph{\{j\textless k\}\textbar F}\{\eta\_j\wedge\eta\emph{k\}\textbar{}}\{L\textsuperscript{2({[}G{]},dg)\}}2.
\tag{G13} {]}

\textbf{Proof.} More generally, let \(E\) be an irreducible unitary
representation of \(K_\infty\) of dimension \(r\), and let
\(F(gk)=\rho(k)^{-1}F(g)\) realize a section of the corresponding
homogeneous bundle. For a unit vector \(e\), Schur averaging gives {[}
\int\emph{\{K}\infty\}\textbar{}\langle\rho(k)\textsuperscript{\{-1\}v,e\rangle\textbar{}}2dk=\{\textbar v\textbar\^{}2\over r\}.
{]} Indeed, the average of the rank-one orthogonal projector onto
\(\mathbf Ce\) commutes with \(\rho\); its trace is one, so it is
\(r^{-1}\operatorname{id}_E\). Right invariance of Haar measure and the
quotient measure formula now give {[}
\int\emph{\{{[}G{]}\}\textbar F\_e\textbar\^{}2dg=\{a\_K\over r\}\int}\{X\_K\}\textbar F\textbar\^{}2dV.
{]} Neatness removes finite stabilizer weights; alternatively the same
equality holds with the usual orbifold weights consistently on both
sides. The cotangent representation of \(U(2)\times U(1)\) is
irreducible of dimension two, while its determinant has dimension one.
In the stated coframe, {[}
\{i\over2\}\eta\wedge\bar\eta\wedge\omega=\textbar{}\eta\textbar\^{}2dV,
\qquad \xi\wedge\bar\xi=4\textbar{}\xi\textbar\^{}2dV. {]} These
pointwise identities prove (G11) and (G12). Insert (G12) into (G1) to
obtain (G13). QED.

\textbf{Normalization consequence.} If a cohomological theta cocycle is
written \(\varphi=\varphi_1\otimes e_1+\varphi_2\otimes e_2\), the
scalar theta lift with Schwartz input \(\varphi_1\) is exactly
\(F_{\eta,e_1}\). Thus no unspecified scalar/cohomological comparison
remains once the unit coframe, actual cocycle, Haar factorization and
\(K_f\) are specified. The lemma does not require the entire form to
belong to a single irreducible automorphic representation. The supplied
human paper uses probability measures on adelic class sets (§4.1,
printed p.12); Haan and GQT use their stated Tamagawa normalizations.
These choices must be converted, not silently equated.

\textbf{Check.} Entire proof reread. The Schur argument applies to the
fiber representation, not to an assumed irreducible global
representation. Both pointwise constants were recomputed in a unit
coframe; \(4\) in (G12), \(2\) in (G13), and the factor \(a_K^{-1}\)
agree with T6.1. Finite stabilizers and all connected components are
included explicitly.

\subsubsection{\texorpdfstring{DEG-L.T6.11 --- An exact finite
Hecke-\(L(3/2)\) energy
formula}{DEG-L.T6.11 --- An exact finite Hecke-L(3/2) energy formula}}\label{deg-l.t6.11-an-exact-finite-hecke-l32-energy-formula}

Use the actual morphism \(f:S\to A\), its polarized coordinate forms,
and their zero extensions to \(X_K\). The oscillator decomposition in
van Hoften--Sempliner, Proposition 3.1.7 and §4.1.7--Proposition 4.1.8
(equivalently the cited Liu and BMM results), gives an orthogonal finite
decomposition of these forms into their actual automorphic summands. For
each occurring summand \(\sigma\), choose its rank-one theta
realization, a nonzero source character vector
\(f_\sigma\in\tau_\sigma\), and scalar Schwartz inputs
\(\phi_{j\sigma}\) such that the coefficient of the \(\sigma\)-part of
\(\eta_j\) along the unit cotangent vector is {[}
F\_\{j\sigma\}=\theta(\bar f\_\sigma,\phi\_\{j\sigma\}). {]} Finite sums
of inputs are allowed. The multiplicity-free oscillator decomposition
and the uniqueness of \((W,\pi)\) with the auxiliary source splitting
fixed are exactly the facts used here; a general multiplicity-one
assertion for all automorphic representations is not needed.

To use Haan's notation, take \(\gamma_\sigma\) to be the splitting
character on the \(U(V)\) side. Changing the source-side splitting to
\(\gamma_\sigma^3\) twists \(\tau_\sigma\) by the corresponding
determinant character and preserves the displayed actual theta lift.
This follows from the splitting-change formula, Wu §2.2, equation (2.9),
printed p.5. Thus each summand can be put into Haan's convention without
changing the form. Define {[} \mathcal Q\_\{j\sigma\}=
\{L\_E(3/2,\operatorname{BC}(\tau\emph{\sigma)\otimes\gamma}\sigma\^{}3)
\over L\_F(3,\epsilon\emph{\{E/F\})\}
,\mathcal Z\^{}\#}\sigma(1;\phi\emph{\{j\sigma\},\phi}\{j\sigma\};f\_\sigma,f\_\sigma),
\tag{G14} {]} where \(\mathcal Z^\#\) is Haan's product of normalized
local integrals for pure tensors and its sesquilinear extension to
finite sums. The source vector's norm is included in this expression.
Haan, Theorem 3.4, identifies \(\mathcal Q_{j\sigma}\) with
\(\|F_{j\sigma}\|^2_{L^2([G])}\); in particular the whole expression is
real and nonnegative, without any claim that each displayed factor
separately is positive.

One then has the \textbf{exact} finite formula {[}
E\_f=\{2\over a\_K\}\sum\emph{\{j,\sigma\}\mathcal Q}\{j\sigma\}, \qquad
\deg\emph{H f{[}S{]}\le
\{4\over a\_K\textsuperscript{2\int\emph{S\omega\^{}2\}
\left(\sum}\{j,\sigma\}\mathcal Q\_\{j\sigma\}\right)}2. \tag{G15} {]}
The corresponding equality with the primitive defect is {[} \deg\emph{H
f}}{[}S{]}= \{4\over a\_K\textsuperscript{2\int\emph{S\omega\^{}2\}
\left(\sum}\{j,\sigma\}\mathcal Q\_\{j\sigma\}\right)}2
-\textbar{}\beta\_0\textbar\^{}2. \tag{G16} {]}

\textbf{Proof.} T6.9 supplies actual finite expansions of the zero
extensions. Project them to the oscillator summands in the cited
direct-sum decomposition, and choose one rank-one model for each
summand. The theta image spans that summand; finite linearity of the
Schwartz input absorbs the finite expansion, with the one-dimensional
source vector fixed. Distinct actual summands are orthogonal for the
adelic invariant inner product. Therefore the norm of each scalar
coefficient is the sum of the norms of its summands. Apply Haan's
bilinear formula to each norm and (G11) to the differential form. Zero
extension makes the full geometric norm equal to the norm on \(S\), so
there is no component restriction inequality at this stage. Finally
apply the Hodge-index identity (G10). QED.

\textbf{Check and scope.} Entire proof reread, including zero extension,
representation orthogonality, finite linearity, source twists, and the
constants in (G15)--(G16). Haan's actual theorem is bilinear, so it
applies to Gram entries before finite inputs are combined. Its global
character factor is absolutely convergent and nonzero at \(3/2\)
(§3.6.1, printed p.9). This proves an exact energy expression and a
degree inequality in terms of finitely many Hecke values, rather than an
asserted equality of the degree with one Hecke value. It is not yet an
effective level bound: the finite geometric/theta realization
coefficients and the finite and archimedean Schwartz inputs are actual
additional data whose size must be bounded. Their existence and the
scalar norm conversion are no longer missing hypotheses.

\subsection{9. Rank-two attack: the wedge has a genuinely finite
cuspidal
spectrum}\label{rank-two-attack-the-wedge-has-a-genuinely-finite-cuspidal-spectrum}

\subsubsection{DEG-L.T6.12 --- Product, toric period, and finite Rallis
expansion}\label{deg-l.t6.12-product-toric-period-and-finite-rallis-expansion}

Let \(\Theta_i\) be two genuine integrated holomorphic one-form theta
lifts from characters \(\chi_i\) of \(T_i=U(W_i)\), with compatible Weil
splittings. Both \(W_i\) have signature \((0,1)\) at the active real
place. Put {[}
W=W\_1\perp W\_2,\quad H=U(W),\quad T=T\_1\times T\_2\subset H,\quad \chi=\chi\_1\boxtimes\chi\_2.
{]} Choose the rank-two cocycle \(\varphi\) to be the tensor/wedge
product of the actual rank-one cocycles, under the direct-sum
identification of Weil representations. This definition includes its
normalization. Let \(\mathcal K_\varphi(g,h)\) be the resulting
holomorphic two-form theta kernel. Then {[}
\Theta\emph{1\wedge\Theta\emph{2
=\int}\{{[}T{]}\}\mathcal K}\varphi(-,t),\overline{\chi(t)},dt.
\tag{G17} {]} This is the theta transform of a toric-period functional
on \(H\), rather than the theta lift of a single character of \(H\).

Assume the finite input is fixed by an open compact subgroup of
\(H(\mathbf A_f)\) and the archimedean input is \(K_{H,\infty}\)-finite,
as holds for the polynomial-Gaussian cocycles just specified. Fix a
common neat \(G\)-level \(K_f\). There are finitely many irreducible
cuspidal \textbf{realizations} \(\pi_\lambda\subset L^2([H])\) and
finitely many orthonormal vectors \(b_{\lambda r}\in\pi_\lambda\) for
which {[} \mathcal K\_\varphi(g,h)=
\sum\emph{\{\lambda,r\}\Theta}\varphi(b\_\{\lambda r\})(g)b\_\{\lambda r\}(h),
\qquad
\Theta\emph{\varphi(b)(g)=\int}\{{[}H{]}\}\mathcal K\_\varphi(g,h)\bar b(h),dh.
\tag{G18} {]} The equality is an equality of smooth two-form-valued
functions, and the sum is genuinely finite. Write {[}
p\_\{\lambda r\}=\int\emph{\{{[}T{]}\}b}\{\lambda r\}(t)\overline{\chi(t)},dt.
{]} Thus the wedge is the finite sum
\(\sum_{\lambda,r}p_{\lambda r}\Theta_\varphi(b_{\lambda r})\).

Let \(\phi\) denote the scalar coefficient of \(\varphi\) in the unit
canonical frame from T6.10. With inner products linear in the first
variable, its exact norm on the full Shimura variety is {[}
\int\emph{\{X\_K\}(\Theta\emph{1\wedge\Theta\emph{2)\wedge
\overline{(\Theta_1\wedge\Theta_2)}
=\{8\over a\_K\}\sum}\lambda L(1,\pi}\lambda\times\chi\emph{V)
\sum}\{r,s\}p}\{\lambda r\}\bar p\_\{\lambda s\}
\mathcal Z\^{}*\emph{\{\lambda\}(1/2;\phi,\phi;b}\{\lambda r\},b\_\{\lambda s\}).
\tag{G19} {]} Here \(\chi_V\) is the retained source-side splitting
character in GQT's notation; \(\mathcal Z^*\) is precisely its
normalized doubling form, extended to finite sums of pure tensors. Its
unramified denominators are retained. The standard \(L\)-function has
degree two. It is not in general a Hecke-character \(L\)-function. No
lower-tower, continuous-spectrum, or residual-spectrum term has been
omitted.

\textbf{Proof.} The direct-sum Schrödinger model is the tensor product
of the two models. At a point of \(T\), its theta sum is the product of
the two theta sums; applying the wedge map to the differential-form
factors gives the identity of kernels. Absolute convergence and
compactness of the two torus quotients allow their integrations to be
combined. This proves (G17). Compatibility of the splitting is
substantive: the character on the \(G\) side is the product of the two
rank-one characters, and the character on the \(H\) side is common to
both lines. A source-character twist permits this common choice. See Wu,
Proposition 2.2 and §2.2 equations (2.8)--(2.9), printed pp.4--5; the
corresponding common auxiliary choice is explicit in the human paper
§4.1.7, printed p.15. The direct-sum theta-product argument also appears
in 032, Remark 4.3 and Proposition 4.6.

The form \(W\) is definite at the active real place, so it has no
nonzero rational isotropic vector. Thus \(H\) is anisotropic over \(F\).
The target \(V\) is definite at a nonactive real place, so \(G\) is
anisotropic as well. Their centers are norm-one tori. Hence both adelic
quotients are compact and their automorphic \(L^2\) spectra are discrete
and cuspidal; cuspidality here is vacuous because there is no proper
\(F\)-parabolic. In particular there is no continuous or residual
contribution from an Eisenstein series on \(H\). Use the standard
discrete admissible decomposition of this compact automorphic quotient,
with finite multiplicities.

We give the orthogonality step explicitly, including distinct copies of
an isomorphic representation. For two orthogonal irreducible
subrepresentations \(\mathcal V_1,\mathcal V_2\subset L^2([H])\), the
mixed convergent-range doubling identity gives the inner product of
theta lifts from \(f_i\in\mathcal V_i\) as the value of a mixed doubling
zeta integral. The identity is the kernel/see-saw identity underlying
GQT, Proposition 11.1(ii) and Proposition 11.2, before an irreducible
representation is fixed. For large \(\operatorname{Re}s\), unfold that
integral. Every rational maximal isotropic plane in \(W\oplus(-W)\) is
the graph of an isometry of \(W\): its projection to either copy of
\(W\) is injective, since a kernel vector would be isotropic in \(W\),
and dimensions then make it an isomorphism. There is therefore only the
graph orbit. The unfolded integral is an integral of a Siegel section
against a matrix coefficient {[}
\langle R(h)f\_1,f\_2\rangle\_\{L\^{}2({[}H{]})\} {]} (up to the fixed
splitting characters, which cancel in the usual doubling normalization).
This coefficient is zero because \(\mathcal V_1\) is invariant and
orthogonal to \(\mathcal V_2\). The zeta integral is identically zero
and stays zero on meromorphic continuation. The mixed theta inner
product is zero. This argument works for distinct orthogonal realization
copies as well as inequivalent representations; it does not assume
global multiplicity one.

For each \(h\), the two-form \(\mathcal K_\varphi(-,h)\) belongs to the
fixed finite-dimensional space \(\mathcal F=H^0(X_K,\Omega^2)\). The
\(H\) action commutes with the \(G\) action and the cohomological
differential, so varying \(h\) preserves this property and the chosen
\(G\)-level. The spectral coefficients \(\Theta_\varphi(b)\) also lie in
\(\mathcal F\). By the preceding orthogonality, the nonzero images
belonging to distinct orthogonal irreducible source realizations are
mutually orthogonal nonzero subspaces of \(\mathcal F\). There can
therefore be only finitely many such realizations. Within each one, the
fixed finite level and the finitely many archimedean compact types of
\(\varphi\) select a finite-dimensional subspace by admissibility. All
other spectral coefficients vanish. The \(L^2\) expansion is
consequently finite; both sides are smooth, so the equality holds
everywhere, proving (G18). Applying the toric functional to this finite
equality is legitimate and introduces no Sobolev convergence issue.

For every retained \(\pi_\lambda\), the global theta lift is nonzero,
hence all its local theta lifts are nonzero. In GQT, the dimensions are
\(n=2,m=3\), \(\epsilon_0=0\), and the target Witt index is \(r=0\).
Thus \(2<3\le4\), \(r\le n\), and the lower-tower hypothesis is empty.
Theorem 11.4(ii), printed pp.54--55, applies at \(s_{m,n}=1/2\) and
gives {[}
\langle\Theta\emph{\phi(b}\{\lambda r\}),\Theta\emph{\phi(b}\{\lambda s\})\rangle
=2L(1,\pi\_\lambda\times\chi\_V)
\mathcal Z\^{}*\emph{\{\lambda\}(1/2;\phi,\phi;b}\{\lambda r\},b\_\{\lambda s\}).
{]} The \(L\)-value here is holomorphic; a residue is not required.
Expand the finite squared norm, use the orthogonality already proved,
and apply (G12). This proves (G19). QED.

\textbf{Check.} Entire proof reread. Definite-at-one-place anisotropy,
the graph-orbit argument, orthogonal-copy issue, and the
finite-dimensional-image argument were checked separately. Finite image
alone would not imply finite spectral support: orthogonality and
admissibility are essential. The toric functional is applied only after
finiteness has been proved. The conjugate in the definition of
\(\Theta_\varphi\) agrees with the rendered GQT §11.2, printed p.52; web
extraction omits that overbar. The factor \(8/a_K=4/a_K\times2\) in
(G19) was checked. This proof uses the standard discrete admissible
spectral decomposition and the cited convergent Siegel--Weil/doubling
identity as background theorems, rather than claiming to reprove them.

\subsubsection{\texorpdfstring{DEG-L.T6.13 --- General degree as a
finite standard-\(L(1)\) norm
sum}{DEG-L.T6.13 --- General degree as a finite standard-L(1) norm sum}}\label{deg-l.t6.13-general-degree-as-a-finite-standard-l1-norm-sum}

The preceding statement extends to every actual map \(f\) in the 032
construction in the following precise, non-effective form. There exist
finitely many rank-two skew-Hermitian spaces \(W_\lambda\), irreducible
cuspidal representations \(\pi_\lambda\) of \(U(W_\lambda)\) with
nonzero lift to \(G\), and finite theta input tensors \(q_{jk,\lambda}\)
such that {[} \deg\emph{H f{[}S{]}
=\{4\over a\_K\}\sum\emph{\{j\textless k\}\sum}\lambda
L(1,\pi\emph{\lambda\times\chi}\{V,\lambda\})
\mathcal Z\^{}}\_\{\lambda\}(1/2;q\_\{jk,\lambda\},q\_\{jk,\lambda\}).
\tag{G20} {]} To make this a defined equality, the data include unitary
\(G(\mathbf A)\)-intertwiners when isomorphic target representations
occur through different theta models or in several copies. The input
tensors encode the actual zero-extended wedge forms through these
identifications. They are not normalized to have size one. The formula
retains the local factors and all input coefficients; it does not assert
that these quantities can be recovered or bounded from the source level
alone.

\textbf{Proof.} Expand the zero extension of each \(\eta_j\) by T6.9.
Each \(\eta_j\wedge\eta_k\) is a finite sum of products of rank-one
theta forms. Apply T6.12 separately to every such product, taking a
common \(G\)-level; all relevant rank-two source groups are anisotropic
at the active place. This places all the scalar wedge coefficients in
the sum of finitely many global rank-two theta images. Every nonzero
such image is an irreducible cuspidal representation by Wu, Theorem 5.3,
printed p.15, since the target is anisotropic. Their sum is a
finite-length unitary representation. Decompose it into orthogonal
irreducible actual copies.

For each resulting copy choose one theta image isomorphic to it and a
unitary intertwiner from that image to the chosen copy. Transport the
orthogonal component of every scalar wedge coefficient back by this
isometry. It is a smooth vector of finite compact type in that theta
image, and hence a finite sum of theta integrals by the definition of
the global theta image (equivalently its irreducibility as an admissible
module). Denote the resulting finite input tensor by \(q_{jk,\lambda}\).
Unitary transport and orthogonality express the squared norm of the
scalar wedge as the sum of squared norms of these theta integrals. Apply
GQT Theorem 11.4(ii) to each sum, extending its bilinear identity by
finite linearity, and insert the result into (G13). The factor is
\(2/a_K\) from (G13) times \(2\) from Rallis. QED.

\textbf{Check and practical distinction.} Entire proof reread, including
the finite-length span and the possible repetition of an abstract target
representation. The unitary identifications are explicitly retained;
simply declaring different source groups or source characters orthogonal
would be invalid. In the special case of one common rank-two source and
splitting, (G20) has the concrete toric-period/input Gram coefficients
of (G19). With different source models, the isotypic comparison data are
additional data. The theorem gives a finite spectral identity and
locates the exact analytic factors; it does not give a canonical single
Hecke value, an effective closed formula in level/conductor alone, or a
bound for the integral realization coefficients. For the desired upper
degree estimate, T6.11 is stronger operationally: it already avoids
rank-two comparison data and uses only Hecke-character values at
\(3/2\).

\section{Bounds, Weil projectors, and the integral map lattice}

Date: 21 November 2026. Actor: GPT-6 Astra Pro, subagent
\texttt{/root/lvalue\_bounds}; effort setting not exposed. Run:
DEG\_L\_20261121. This is a working research contribution to the root
paper, not a claim to solve T6 or T7. Every assertion below is either
explicitly conditional, an elementary derived finding with proof, or an
identified cited background result. All degrees are geometric
intersection degrees, without factorial normalization.

\subsection{Executive finding}\label{executive-finding}

The strongest positive reduction obtained here is \textbf{B8}: in a
polarized family of abelian varieties of Weil type of dimension \(2n\),
a uniformly bounded-degree effective \(n\)-cycle with \emph{any nonzero
Weil component} supplies a bounded-division representative of every
fixed flat rational Weil class. The multiplier and both signed degrees
are \(O(D^2)\), with explicit constants depending only on the fixed
family, polarization type, field action, and target class. No varying
full CM field or varying Hecke projector is required. The proof gives a
fixed polynomial correspondence in the generic imaginary-quadratic
action, a comass estimate on invariant harmonic forms, and an integral
two-dimensional Cramer calculation. This removes the denominator and
isolation gates once bounded geometric witnesses exist; the existence of
those witnesses in the open families remains an actual hypothesis.

The analytic findings are complementary. The genuine \(U(1)\to U(3)\)
Rallis formula uses \(L_E(3/2,\chi)\), not \(L_E(1,\chi)\); its global
factor has elementary conductor-independent upper/lower bounds for fixed
field degree. It does not compute the degree of the unnormalized 032
image surface. Edge-of-strip bounds under the admitted Alpöge input and
conditional Masullo input apply only after identifying the correct
normalized \(L\)-function, special point, local vectors, measure, and CM
extraction coefficients.

\subsection{Source findings and exact
scope}\label{source-findings-and-exact-scope}

\textbf{S032.} OpenAI family 032, \emph{The rational Hodge conjecture
for CM abelian varieties}, §4.4 Lemma 4.4 (printed pp.20--21) fixes
finite Schwartz data and permits all sufficiently deep principal levels;
§7 proof of Theorem 2.5, printed p.46, explicitly states that pulling to
a deeper principal cover multiplies the mixed integral by the cover
degree. The theorem then selects a nonzero term and rational Hodge maps
by span and finite expansion; it supplies no estimate for coefficients
of those maps. These are statements used and rechecked in this run, not
independent reproofs of 032.

\textbf{S-Rallis1.} Jaeho Haan, \emph{The Bessel Period of U(3) and U(2)
involving a non-tempered representation}, arXiv:1403.5061v7, §3.3,
Theorem 3.4, printed p.8, gives, in that paper's Tamagawa and splitting
conventions, {[}
\langle\theta(\bar f\_1,\varphi\emph{1),\theta(\bar f\_2,\varphi\emph{2)\rangle
=
\frac{L_E(3/2,\operatorname{BC}(\tau)\otimes\gamma^3)}{L_F(3,\epsilon_{E/F})}
\prod\emph{v Z\_v\^{}\{\#\}(1,f}\{1,v\},f}\{2,v\},\Phi}\{1,v\}). {]}
Here \(\tau\) is an automorphic character of \(U(1)\) and \(\gamma\) is
the chosen unitary splitting character; §3.1 (3.1)--(3.2) defines the
normalization of local zeta integrals. §3.6.1, printed p.9, explicitly
records absolute convergence and nonvanishing of the numerator at
\(3/2\). This is a formula for integrated theta lifts with source
character, and cannot be substituted for an arbitrary theta-kernel
evaluation or wedge in 032 without comparison.

URL: https://arxiv.org/pdf/1403.5061 . Retrieved primary ref:
\texttt{turn140259939view0}; exact locator find:
\texttt{turn301895159view3}.

\textbf{S-Rallis2.} Gan--Qiu--Takeda, \emph{The Regularized Siegel--Weil
Formula (The Second Term Identity) and the Rallis Inner Product
Formula}, arXiv:1207.4709v3, Theorem 11.4, printed pp.54--55, under its
explicit range and cuspidality/first-occurrence hypotheses, gives the
unitary point \((m-n+1)/2\), with twisting character and normalized
doubling integrals. Its range \(n<m\leq2n\) excludes \(U(1)\to U(3)\).
The \(U(2)\to U(3)\) point is \(1\), but the standard degree-two
\(L\)-function is not thereby a Hecke-character \(L\)-function. A source
spectral decomposition or endoscopic factorization must be supplied. The
theorem is used as background and its range checked, not reproved.

URL: https://arxiv.org/pdf/1207.4709 . Refs:
\texttt{turn640628982view1}, \texttt{turn819750783view0}.

\textbf{S-Alpöge.} The user's Dirichlet GRH input is admitted. Exact
supplied Lean declarations: \texttt{ALPOGE\_Main.lean},
\texttt{LFunction\_eq\_zero\_re\_of\_isPrimitive\ :\ PrimitiveDirichletLZerosOnLine},
and
\texttt{LFunction\_ne\_zero\_of\_half\_lt\_re\ :\ DirichletLNonvanishingRightOfHalf};
definitions in \texttt{ALPOGE\_Statements.lean}. The supplied
\emph{Remark}, step 20, printed p.4, announces GRH; steps 9, 12, 18
concern lift norms, Rallis, and indices respectively. The full imported
archive is not supplied, and this run did not compile it. No claimed new
divisibility theorem is inferred from those steps or from the mapping
file.

\textbf{S-Masullo.} The full draft is not supplied. The prompt
attributes Hecke GRH to §1.2. The supplied
\texttt{MASULLO\_STATEMENTS\_INDEX.md} is an index of openings, not
theorem statements or proofs. It records Chapter 12, Theorem 12.5.1.1,
PDF p.519, and Chapter 13, Theorem 13.3.1.1, PDF p.522, as trace
formulas, and records an unfinished-looking opening labelled 13.3.1.2
(at the text-layer level). None is silently promoted to a precise
theorem of Hecke GRH in this report. All use of Hecke GRH is conditional
on the claim specified by the user, not certified from the index.

\subsection{B1. Elementary uniform bounds in the absolutely convergent
half-plane}\label{b1.-elementary-uniform-bounds-in-the-absolutely-convergent-half-plane}

\textbf{Statement.} Let \(K\) have degree \(d\), and let \(\chi\) be a
unitary Hecke character. For real \(\sigma>1\), using the finite Euler
product with ramified factors equal to one, {[}
\left(\frac{\zeta(2\sigma)}{\zeta(\sigma)}\right)\^{}d
\leq \textbar L\_K(\sigma,\chi)\textbar{}\leq \zeta(\sigma)\^{}d. {]}
More sharply,
\(\zeta_K(2\sigma)/\zeta_K(\sigma)\leq|L_K(\sigma,\chi)|\leq\zeta_K(\sigma)\).
No RH or GRH is needed.

\textbf{Proof.} For every unramified prime ideal \(\mathfrak p\),
unitarity gives \(|\chi(\mathfrak p)|=1\). Thus {[}
(1+N\mathfrak p\textsuperscript{\{-\sigma\})}\{-1\}\leq
\textbar1-\chi(\mathfrak p)N\mathfrak p\textsuperscript{\{-\sigma\}\textbar{}}\{-1\}
\leq(1-N\mathfrak p\textsuperscript{\{-\sigma\})}\{-1\}. {]} The same
inequalities hold with the middle term replaced by one at ramified
primes. Absolute convergence permits multiplication. The left product is
\(\zeta_K(2\sigma)/\zeta_K(\sigma)\). Above each rational prime there
are at most \(d\) prime ideals, each with norm at least that prime.
Comparing each factor gives the stated bounds by rational zeta
functions. All factors are positive and the products converge. QED.

\textbf{Application.} In S-Rallis1 with \([F:\mathbb Q]=d\) and
\([E:\mathbb Q]=2d\), the absolute value of the global ratio lies
between {[} \frac{\zeta(3)^d}{\zeta(3/2)^{2d}} \quad\hbox{and}\quad
\zeta(3/2)\textsuperscript{\{2d\}\left(\frac{\zeta(3)}{\zeta(6)}\right)}d.
{]} This follows by applying B1 separately to numerator and denominator.
It controls no local \(Z_v^\#\) and no change of measure, input scale,
or cycle normalization.

\subsection{B2. Exact edge bounds and what each GRH input
buys}\label{b2.-exact-edge-bounds-and-what-each-grh-input-buys}

Lamzouri--Li--Soundararajan, \emph{Conditional bounds for the least
quadratic non-residue and related problems}, arXiv:1309.3595v3, Theorem
1.5 (printed pp.3--4), gives for primitive Dirichlet characters of
conductor \(q\geq10^{10}\) under GRH, with {[}
A(q)=\log\log q-\log2+\tfrac12+1/\log\log q, {]} {[}
\textbar L(1,\chi)\textbar{}\leq 2e\^{}\gamma A(q),\qquad
\textbar L(1,\chi)\textbar\^{}\{-1\}\leq\frac{12e^\gamma}{\pi^2}
\left(A(q)+\frac{14\log\log q}{\log q}\right). {]} The constant in the
second inequality is \(12e^\gamma/\pi^2\), not \(12e^{2\gamma}/\pi^2\).
An older hosted text extraction returned the latter; the latest arXiv
version and Lumley's Theorem 1.1 provide the correct former value. Small
conductors require their own finite verification or a separate bound.
Under the user's admitted Alpöge theorem, this consequence no longer has
a GRH condition. URL: https://arxiv.org/pdf/1309.3595 ; ref
\texttt{turn301895159view0}.

For primitive nontrivial unitary Hecke characters whose completed
\(L\)-function is entire, write the completed function in degree
\(d=[K:\mathbb Q]\) over \(\mathbb Q\), with conductor
\(q(\chi)=|D_K|N\mathfrak f_\chi\) and archimedean parameters
\(\kappa_j\). Allysa Lumley, \emph{Explicit bounds for L-functions on
the edge of the critical strip}, arXiv:1804.09850v1, §1.1 and Theorems
2.1--2.2 (printed pp.3--4), applies to {[}
C(\chi)=\frac{q(\chi)}{\pi^d}\prod\_\{j=1\}\^{}d\left\textbar{}\frac{1+\kappa_j}{2}\right\textbar.
{]} For fixed degree, its bounds have orders \((\log\log C)^{d}\) and
\((\log\log C)^{-d}\); its explicit formulas apply when
\(\log C\geq23d\). Unit modulus of unramified Hecke parameters supplies
the required Ramanujan hypothesis; the remaining GRH input is
conditional on Masullo for general \(K\). Pole cases (including the
trivial character) are excluded. Field degree and archimedean type
cannot be hidden while allowing them to vary. URL:
https://arxiv.org/pdf/1804.09850 ; ref \texttt{turn347054114view0}.

These are known results applied, not new theorems. In particular, this
report does not assert an edge bound for a nonabelian \(U(2)\) standard
\(L\)-function from Hecke GRH. Nor does GRH give a positive lower bound
at the center: the functional equation can force a central zero when a
self-dual root number is \(-1\).

\subsection{B3. Level, vector scale, and arithmetic conductor are
distinct}\label{b3.-level-vector-scale-and-arithmetic-conductor-are-distinct}

\textbf{Statement.} If \(\pi:Y\to X\) is a finite étale cover of degree
\(d\) and \(\eta\) is a top-degree form on \(X\), then
\(\int_Y\pi^*\eta=d\int_X\eta\). Consequently a nonzero pulled-back
theta form has unnormalized Petersson norm squared multiplied by \(d\).
If its underlying automorphic representation, infinity type, and finite
vector are fixed, the representation's arithmetic and analytic
conductors stay fixed.

\textbf{Proof.} Integration over evenly covered coordinate sets and a
partition of unity gives the integral identity. Apply it to the positive
top form defining the norm. Conductors are invariants of the fixed
representation and archimedean parameters, not of which smaller compact
open subgroup happens to fix its vector. QED.

This is also explicitly recorded for the mixed period in 032 §7 p.46. It
is an exact level-growth statement for a chosen presentation, not a
lower bound for the intrinsic minimum complexity \(c(s)\) in Route 2.
Similarly scaling a Schwartz input by \(u\) scales a theta lift by \(u\)
and its norm squared by \(|u|^2\), while leaving its representation's
\(L\)-function unchanged. The local integrals and/or input pairing carry
that factor.

\subsection{B4. Polarized isogenies: rational degree can stay fixed
while denominators
grow}\label{b4.-polarized-isogenies-rational-degree-can-stay-fixed-while-denominators-grow}

\textbf{Statement.} Let \(\psi:A\to B\) be an isogeny of
\(g\)-dimensional polarized abelian varieties. Suppose
\(\psi^*c_1(L_B)=c\,c_1(L_A)\) and \(L_A,L_B\) have the same
polarization degree. For a codimension-\(p\) cycle \(Z\), {[}
\deg\psi=c\^{}g,\qquad \deg\_\{L\_A\}(\psi\^{}*Z)=c\^{}p\deg\_\{L\_B\}Z.
{]}

\textbf{Proof.} Integrating the \(g\)th power of the polarization
identity gives the first equality. By projection formula, {[}
\int\_A\psi\^{}*{[}Z{]}c\_1(L\_A)\^{}\{g-p\}
=c\^{}\{-(g-p)\}\deg\psi\int\_B{[}Z{]}c\_1(L\_B)\^{}\{g-p\}, {]} which
is the second. QED.

The rational cycle \(c^{-p}\psi^*Z\) has unchanged degree but its
displayed coefficient has denominator \(c^p\) when \(c\) is an integer.
For multiplication by \(m\) on a fixed variety, \(c=m^2\). This is not a
proof that all representatives have large denominators: it is a precise
diagnosis of the natural pullback representative.

\subsection{B5. A rigorous norm-to-degree bridge for surface
maps}\label{b5.-a-rigorous-norm-to-degree-bridge-for-surface-maps}

\textbf{Statement.} Let \(Y\) be a smooth projective surface, \(\omega\)
a Kähler class, \(V=\int_Y\omega^2\), and \(f:Y\to A\) a morphism with
two-dimensional image \(Z\) and generic degree \(\delta\). Choose a
translation-invariant representative \(\omega_A\in c_1(L)\) of a
polarization, and invariant holomorphic forms
\(\alpha_1,\ldots,\alpha_g\) such that
\(\omega_A=i\sum_j\alpha_j\wedge\bar\alpha_j\). Put {[} E\_f=\int\_Y
f\^{}*\omega\_A\wedge\omega
=\sum\_j\textbar f\^{}*\alpha\emph{j\textbar{}}\omega\^{}2,\qquad
\textbar u\textbar\_\omega\^{}2=i\int\_Yu\wedge\bar u\wedge\omega. {]}
Then {[} \deg\_L Z=\frac{\int_Y(f^*\omega_A)^2}{\delta}
\leq\frac{E_f^2}{\delta V}. {]}

\textbf{Proof.} Proper pushforward gives \(f_*[Y]=\delta[Z]\) and
projection formula gives the equality. With \(\beta=f^*\omega_A\),
decompose \(\beta=(E_f/V)\omega+\beta_0\), where
\(\int\beta_0\omega=0\). The Hodge index theorem on the surface gives
\(\int\beta_0^2\leq0\), so \(\int\beta^2\leq E_f^2/V\). The displayed
energy identity follows by linearity from the specified polarization
representative. QED.

This is a direct application of the classical surface Hodge index
theorem in a form adapted to the present problem, not a claim to reprove
that theorem. Under a cover of degree \(d\), \(E_f\) and \(V\) both
scale by \(d\). Thus the bound \(E_f^2/V\) for the degree of the
pushforward class scales by \(d\). For the reduced image, the generic
mapping degree \(\delta\) also scales by \(d\), and the displayed bound
remains unchanged; this matches B3 exactly.

For the sufficiently deep principal ball-quotient surfaces used in 032,
the actual pullbacks \(f^*\alpha_j\) do have finite expansions
\(\sum_a b_{ja}\theta_a\) in genuine integrated theta lifts from
characters of one-dimensional unitary groups. The existence of these
expansions is supplied by Bergeron--Millson--Mœglin, \emph{The Hodge
conjecture and arithmetic quotients of complex balls}, Acta Math.216
(2016), Theorem 9.1, pp.73--75, with the range in §9.1, p.72: for
\(m=3\) and \((a,b)=(0,1)\), its condition is \(4<6\). The supplied van
Hoften--Sempliner paper makes the all-\(H^{1,0}\) spanning conclusion
explicit in the proof of Proposition 4.1.8, pp.15--16; Proposition
3.1.7, p.11, identifies the oscillator constituents. The published BMM
proof and the supplied Proposition 4.1.8 proof were read. This is
imported classification/spanning, not a new proof of those theorems.

For completeness, only finitely many direct-limit theta vectors are
needed for each fixed form. Put them at a common finer level and average
over the desired compact open subgroup. The original form is invariant,
and equivariance moves the finite average to the Schwartz inputs, so the
expansion exists at that level. Equality in \(H^{1,0}\) is equality of
holomorphic forms by the injectivity into Betti cohomology for a compact
Kähler variety. If necessary one first deepens the principal level so
the arithmetic \(U\) and \(SU\) groups agree: their determinant values
are norm-one algebraic units, hence roots of unity, and sufficiently
deep congruence kills the finite set of nontrivial roots. T6.9 gives the
full argument and its scope.

With norm squares \(Q_a\) computed by a valid rank-one Rallis formula,
and with the scalar/geometric metric and measure conversions included in
\(Q_a\), the triangle inequality gives the concrete bound {[} \deg\_L
Z\leq\frac{1}{\delta V}
\left\{\sum\_j\left(\sum\emph{a\textbar b}\{ja\}\textbar{}\sqrt{Q_a}\right)\textsuperscript{2\right\}}2.
{]} The existence gate is therefore removed. What remains is a
quantitative bound on the coefficients \(b_{ja}\) relative to the
integral Hom lattice, together with CM periods, test-vector factors and
metric/measure normalization. Orthonormalizing the theta basis does not
control those integral coefficients. B10 below reduces them to an
integral energy-lattice covolume estimate, and B11 shows exactly why
that estimate cannot be inferred from finite generation or
orthonormalization alone.

\subsection{B6. Fixed generic-field projector onto Weil
classes}\label{b6.-fixed-generic-field-projector-onto-weil-classes}

\textbf{Setup.} Let \(A\) be a polarized abelian variety of dimension
\(2n\), carrying an action of an order \(\mathcal O\subset K\) in an
imaginary quadratic field. Assume the Rosati involution restricts to
conjugation on \(K\) and the two eigenspaces in \(H^{1,0}\) both have
dimension \(n\) (Weil type). Let {[}
W=\bigwedge\_K\textsuperscript{\{2n\}H}1(A,\mathbb Q)\subset H\^{}\{2n\}(A,\mathbb Q).
{]} Here the inclusion means the rational rank-two subspace whose
complexification is the sum of the two extreme \(K\)-eigenlines. Choose
a nonzero \(a\in\mathcal O\) such that \(a/\bar a\) is not a root of
unity, and let \(T=[a]^*\) on cohomology. Such an \(a\) exists: only
finitely many arguments arise from roots of unity in \(K\), whereas the
arguments of nonreal elements \(M+\tau\) for a fixed nonreal
\(\tau\in\mathcal O\) vary with integers \(M\) and tend to zero.

\textbf{Statement.} There is a polynomial \(P(X)\in\mathbb Z[X]\) and a
positive integer \(R\), depending only on \((K,n,a)\), such that
\(e=P(T)/R\) is the projector onto \(W\) on \(H^{2n}\). It is induced by
a fixed rational algebraic correspondence and is self-adjoint for the
Poincaré pairing on that degree. Its denominator is independent of the
CM point or Hecke level. One possible choice is {[}
\lambda\emph{r=a\textsuperscript{r\bar a}\{2n-r\}~(0\leq r\leq2n),\quad
f\_W(X)=(X-a\textsuperscript{\{2n\})(X-\bar a}\{2n\}),\quad
f\_R(X)=\prod}\{r=1\}\^{}\{2n-1\}(X-\lambda\_r),\quad
R=\operatorname{Res}(f\_W,f\_R)\textgreater0, {]} with \(P=vf_R\) for
integral polynomials satisfying \(u f_W+v f_R=R\).

\textbf{Proof.} Decomposing \(H^1_\mathbb C=V_a\oplus V_{\bar a}\), each
summand has dimension \(2n\), and {[}
H\textsuperscript{\{2n\}\emph{\mathbb C=\bigoplus}\{r=0\}}\{2n\}\bigwedge\textsuperscript{rV\_a\otimes\bigwedge}\{2n-r\}V\_\{\bar a\}.
{]} The corresponding eigenvalues are \(\lambda_r\). They are distinct
since \(a/\bar a\) has infinite order. Each polynomial displayed is
fixed by conjugation and has algebraic-integer coefficients, hence lies
in \(\mathbb Z[X]\); the two factors are coprime. The Sylvester matrix
for these monic polynomials has determinant their resultant. Multiplying
that matrix by its integral adjugate and extracting the column
corresponding to the constant polynomial gives integral \(u,v\) with
\(u f_W+v f_R=R\) (a simultaneous sign change adjusts the determinant
convention). Thus \(P/R\) takes value one at the two extreme eigenvalues
and zero at every other eigenvalue. Since \(T\) is diagonalizable, it is
the asserted projector.

Every power \(T^j\) is induced by pullback along the algebraic isogeny
\([a^j]\), so the polynomial defines an algebraic correspondence. Its
cohomological projector property is what is asserted: \textbf{no
Chow-group identity \(e^2=e\) is claimed} merely from this calculation.
For a full homological-motive statement, compose it with the algebraic
degree projector \(\pi_{2n}=\prod_{j\ne2n}([2]^*-2^j)/(2^{2n}-2^j)\),
where \(0\leq j\leq4n\); this acts as the identity on \(H^{2n}\) and as
zero in all other degrees. Thus it defines an idempotent on \(h^{2n}\)
in homological motives, and on numerical cycle classes, without claiming
a Chow idempotence relation.

To prove self-adjointness on this degree, the Poincaré pairing of two
\(T\)-eigenvectors can be nonzero only if
\(\lambda_r\lambda_s=\deg[a]=N(a)^{2n}\): this follows by applying
\([a]^*\) to their product in top degree. Infinite order of \(a/\bar a\)
makes this condition exactly \(r+s=2n\). Consequently the two extreme
eigenlines pair only with each other and are orthogonal to all the
intervening eigenspaces. Projection onto their sum is
Poincaré-self-adjoint. Finally,
\(R=f_R(a^{2n})\overline{f_R(a^{2n})}>0\). QED.

A useful explicit expression is {[}
R=N(a)\^{}\{n(2n-1)\}\prod\_\{k=1\}\textsuperscript{\{2n-1\}\textbar a}k-\bar a\textsuperscript{k\textbar{}}2.
{]} It follows by factoring
\(f_R(a^{2n})=a^{n(2n-1)}\prod_{k=1}^{2n-1}(a^k-\bar a^k)\) and taking
the field norm.

If \(P=\sum_jb_jX^j\), put {[}
B\_P=\sum\_j\textbar b\_j\textbar N(a)\^{}\{jn\}. {]} For effective
codimension-\(n\) cycles \(Z\) of degree \(D\), \(P([a]^*)Z\) is an
integral signed cycle whose positive and negative parts both have degree
at most \(B_PD\). Indeed \([a]^*c_1(L)=N(a)c_1(L)\) by the Rosati
condition, so B4 gives degree multiplier \(N(a)^n\) for pullback on
codimension \(n\); use it term by term. This works fibrewise in any
family with the fixed action and polarization data.

\subsection{B7. Uniform Weil-component bound from effective
degree}\label{b7.-uniform-weil-component-bound-from-effective-degree}

\textbf{Statement.} In B6 put \(\omega=c_1(L)\),
\(V_A=\int_A\omega^{2n}/(2n)!=\chi(L)\), and {[}
q(u,v)=(-1)\^{}n\int\emph{Au\wedge v\quad(u,v\in W}\mathbb R). {]} Then
\(q\) is positive definite. For every effective \(n\)-cycle \(Z\) of
degree \(D=\int_Z\omega^n\), {[}
\textbar e{[}Z{]}\textbar\_q\leq\frac{D}{n!\sqrt{V_A}}. {]} The constant
depends only on polarization type and \(n\), hence is uniform on the
whole Weil family, without a compactness restriction on its base.

\textbf{Proof.} The Rosati identity makes \(\omega\) a mixed-eigenvalue
class in \(V_a\otimes V_{\bar a}\): its \([a]^*\) eigenvalue is \(N(a)\)
on degree two, different from \(a^2\) and \(\bar a^2\). Wedge product by
\(\omega\) kills each extreme degree-\(2n\) eigenline because it would
require \(2n+1\) vectors in one \(2n\)-dimensional eigenspace. Thus
\(W\) is primitive. Weil type makes both lines of Hodge type \((n,n)\).
The Hodge--Riemann bilinear relations give positivity of \(q\); the
Hodge star on these primitive middle-degree real classes is
\(*=(-1)^n\).

Give \(A\) its flat translation-invariant metric associated with
\(\omega\). Every harmonic form is translation invariant, so for its
harmonic representative \(\eta\in W_\mathbb R\), the pointwise Euclidean
norm is constant and equals \(\|\eta\|_q/\sqrt{V_A}\). Its comass is at
most its Euclidean norm. The complex analytic cycle \(Z\) has volume
\(D/n!\) by Wirtinger's identity, with multiplicities; its singular
locus has measure zero. Therefore {[}
\left\textbar{}\int\_Z\eta\right\textbar{}\leq\frac{D}{n!\sqrt{V_A}}\textbar{}\eta\textbar\_q.
{]} By B6, \(e\) is self-adjoint for Poincaré pairing, so for
\(w=e[Z]\), {[}
\textbar w\textbar\_q\^{}2=\textbar q(w,w)\textbar=\left\textbar{}\int\_A{[}Z{]}\wedge w\right\textbar=\left\textbar{}\int\_Zw\right\textbar.
{]} If \(w=0\) the result holds; otherwise apply the preceding
inequality with \(\eta=w\) and divide by \(\|w\|_q\). QED.

\subsection{B8. Effective witness implies bounded division, with
quadratic
bounds}\label{b8.-effective-witness-implies-bounded-division-with-quadratic-bounds}

\textbf{Statement.} Keep a connected polarized Weil-type family as in
B6, with a globally defined \(\mathcal O\)-action (after finite base
change if necessary), fixed \(n,K,a,L\), and fixed rational flat target
\(t\in W\). Put {[}
\Lambda=W\cap H\^{}\{2n\}(A,\mathbb Z),\qquad \delta\_\Lambda=\det(q(e\_i,e\_j))\textgreater0
{]} for any integral basis of this rank-two saturated lattice. Choose
\(b\geq1\) with \(\tau=bt\in\Lambda\), and write \(h=N(a)^n\). The
constants \(V_A\), \(\delta_\Lambda\), and \(\|\tau\|_q\) are locally
constant along the family and so fixed on a connected base.

Suppose a fibre contains an effective \(n\)-cycle \(Z\) of degree \(D\)
with \(e[Z]\ne0\). Then there are an integer \(N\geq1\) and effective
\(Z^+,Z^-\) with {[} Nt={[}Z\^{}+{]}-{[}Z\^{}-{]}, {]} where, with
\(M=RD/(n!\sqrt{V_A})\), {[}
N\leq\frac{b h M^2}{\sqrt{\delta_\Lambda}},\qquad
\deg Z\^{}\pm\leq\frac{2h\|\tau\|_q M B_P D}{\sqrt{\delta_\Lambda}}. {]}
Thus both are \(O(D^2)\), uniformly in the fibre. In particular a
uniform bound for degrees of nonzero Weil witnesses on a Zariski-dense
set gives bounded division on that set. The proper parameter-space
argument of Route 2 then applies, once supplied separately with its
stated hypotheses.

\textbf{Proof.} Let \(V_Z=P([a]^*)Z\) as an integral signed cycle and
\(v=[V_Z]=R e[Z]\). This is a nonzero element of \(\Lambda\), with
\(\|v\|_q\leq M\) by B7. The characteristic polynomial of \(T\) on \(W\)
is \(f_W\). Since \(a/\bar a\) is not a root of unity, \(a^{2n}\) is
nonreal, so \(f_W\) is irreducible over \(\mathbb Q\). Hence \(v,Tv\)
are a rational basis of \(W\): otherwise \(v\) would be a rational
eigenvector of \(T\) with a rational eigenvalue.

Use any oriented integral basis of \(\Lambda\). Its determinant {[}
\Delta=\det(v,Tv)\in\mathbb Z\setminus\{0\} {]} therefore satisfies
\(\Delta\tau=m v+\ell Tv\), where the integral Cramer numerators are
\(m=\det(\tau,Tv)\) and \(\ell=\det(v,\tau)\). On \(W\), \(T\) is a
similitude of \(q\) of multiplier \(\deg[a]=N(a)^{2n}=h^2\), so
\(\|Tv\|_q=h\|v\|_q\). The area of a parallelogram in this metric is
\(\sqrt{\delta_\Lambda}\) times its lattice determinant. The elementary
area bound consequently gives {[}
\textbar{}\Delta\textbar{}\leq hM\^{}2/\sqrt{\delta_\Lambda},\quad
\textbar m\textbar{}\leq\textbar{}\tau\textbar\_q
hM/\sqrt{\delta_\Lambda},\quad
\textbar{}\ell\textbar{}\leq\textbar{}\tau\textbar\_q
M/\sqrt{\delta_\Lambda}. {]} Let \(N=b|\Delta|\), and let
\(C=\operatorname{sgn}(\Delta)(mV_Z+\ell[a]^*V_Z)\). Then \([C]=Nt\). By
the final part of B6 the positive and negative degrees of \(V_Z\) are at
most \(B_PD\), and those of \([a]^*V_Z\) are at most \(hB_PD\).
Splitting the signed cycle \(C\) into positive and negative parts thus
bounds each by {[}
(\textbar m\textbar+h\textbar{}\ell\textbar)B\_PD\leq 2h\textbar{}\tau\textbar\_q
M B\_PD/\sqrt{\delta_\Lambda}. {]} This proves the claim. The
intersection lattice and its pairing vary as an integral local system,
so its discriminant and the norm of a fixed flat target are constant;
\(V_A\) is fixed by polarization type. No choice of a global basis is
needed to make the constants uniform. QED.

\textbf{Stronger exact index identity.} If \(\mu=a^{2n}\) under the
fixed complex embedding, then {[}
\textbar{}\Delta\textbar=\frac{|\operatorname{Im}\mu|}{\sqrt{\delta_\Lambda}},\textbar v\textbar\_q\^{}2.
{]} Indeed \(T\) is multiplication by the nonreal scalar \(\mu\) on the
real plane \(W\), its characteristic polynomial has trace
\(2\operatorname{Re}\mu\), and it is a positive-orientation
\(q\)-similitude. In a \(q\)-orthonormal basis it is the
rotation-dilation matrix for \(\mu\) or \(\bar\mu\); the area of
\((v,Tv)\) is \(|\operatorname{Im}\mu|\|v\|^2\). Dividing by the lattice
covolume proves the formula. This makes the last arithmetic index
precise, without claiming it is automatically an automorphic
\(L\)-value.

\textbf{The quadratic size is sharp for extraction from one lattice
vector.} In the positive lattice \(\Lambda=\mathbb Z[i]\), let \(v=m+i\)
and target \(\tau=1\). If \(N\tau\) belongs to the principal ideal
\(v\mathbb Z[i]\), then {[}
\frac{N}{m+i}=\frac{Nm}{m^2+1}-i\frac{N}{m^2+1}\in\mathbb Z{[}i{]}, {]}
so \(m^2+1\mid N\). Since \(\|v\|^2=m^2+1\), extraction using \(v\) and
any fixed integral multiplication operator can require an integer of
quadratic size in \(\|v\|\). This is an arithmetic sharpness statement
for that extraction method. It does not establish growth of intrinsic
cycle complexity, since other geometric cycles might represent the
target more efficiently.

\textbf{What remains.} The theorem does not create \(Z\) in the non-CM
fibres, and 032 gives neither a uniform degree bound for such \(Z\) on a
dense orbit nor a deformation of one actual full cycle class. It proves
that a separate unbounded full-CM-field projector is unnecessary in the
Weil-family route. Once geometric witness degrees are bounded, the
two-dimensional integral extraction is automatically bounded.

\subsection{B9. A precise failure of archimedean norm control for
arbitrary
projectors}\label{b9.-a-precise-failure-of-archimedean-norm-control-for-arbitrary-projectors}

On the lattice \(\mathbb Z^2\) with Euclidean pairing, the orthogonal
projector to \(\mathbb Q(m,1)\) is {[} E\_m=\frac1{m^2+1}

\begin{pmatrix}m^2&m\\m&1\end{pmatrix}

. {]} It has operator norm one, while its exact common denominator is
\(m^2+1\), because its lower-right entry is \(1/(m^2+1)\). The
calculation follows by the rank-one projector formula \(vv^t/(v^tv)\).
This example already lives in
\(\operatorname{End}(\mathbf1\oplus\mathbf1)\) in homological or
numerical motives: all matrices are algebraic correspondences between
two copies of a point. It refutes any abstract inference from bounded
archimedean projector norm to bounded integral denominator. It does not
refute B6, which fixes the generic-field projector, or the existence of
other bounded representatives in a geometric family.

Similarly \(x_m=(1,1/m)\) has norm in \([1,\sqrt2]\) while its pairing
with \((0,1)\) is \(1/m\). Lower bounds for two individual norms
therefore give no lower bound for the mixed period needed in 032. The
positive continuous comparison used by 032 is a separate input, whose
quantitative approximation costs must be controlled.

\subsection{B10. An integral Hom covolume bound replaces individual
coefficient
bounds}\label{b10.-an-integral-hom-covolume-bound-replaces-individual-coefficient-bounds}

\textbf{Setup.} Let \(Y\) be a smooth projective variety of dimension
\(d\ge1\), with an integral ample class \(H\), and fix a base point for
its Albanese map \(\alpha:Y\to B=\operatorname{Alb}(Y)\). Let
\((A_i,L_i)\) be polarized abelian varieties, for \(1\le i\le k\). Put
{[} \Lambda\_i=\operatorname{Hom}(B,A\_i),\qquad E\_i(u)=\int\_Y
(u\alpha)\^{}*c\_1(L\_i)H\^{}\{d-1\},\qquad V=\int\_YH\^{}d. {]} Here
maps are homomorphisms; translations do not affect the forms, energies
or degrees. Assume \(r_i=\operatorname{rank}\Lambda_i>0\). Let \(q_i\)
be the real quadratic metric with square norm \(E_i\), let \(C_i\) be
the Euclidean covolume of \(\Lambda_i\) for \(q_i\), and let \(\mu_i>0\)
be any known lower bound for \(E_i(u)\) for nonzero integral \(u\). One
may always take \(\mu_i=1\). Set {[}
v\_r=\frac{\pi^{r/2}}{\Gamma(1+r/2)},\qquad
R\_i=\frac{2^{r_i}}{v_{r_i}}\frac{C_i}{\mu_i^{(r_i-1)/2}}. {]}

\textbf{Statement.} The lattice \(\Lambda_i\) contains \(r_i\)
independent homomorphisms \(u_{ij}\), spanning
\(\Lambda_i\otimes\mathbb Q\), such that \(E_i(u_{ij})\le R_i^2\).
Suppose a specified real-multilinear period \(P\) on
\(\prod_i(\Lambda_i\otimes\mathbb R)\) is nonzero. Some tuple of these
homomorphisms has \(P(u_{1j_1},\ldots,u_{kj_k})\ne0\). If the period is
the integral of a wedge of \(2d\) pulled-back one-forms, its
nonvanishing implies that the product map {[}
F=(u\_\{1j\_1\}\alpha,\ldots,u\_\{kj\_k\}\alpha):Y\longrightarrow\prod\_i
A\_i {]} has a \(d\)-dimensional image \(Z\). For the product
polarization \(L=\bigotimes_i\operatorname{pr}_i^*L_i\), and the actual
generic degree \(\delta_F\) of \(F\) onto \(Z\), {[} \deg\_L
Z\le\frac{(\sum_iR_i^2)^d}{\delta_F V^{d-1}}. \tag{B10.1} {]} For 032,
\(d=2\) and \(k=4\). More generally, if a target polarization \(L'\)
satisfies \(B c_1(L)-c_1(L')\) nef for a known \(B>0\), the right side
is multiplied by \(B^d\).

\textbf{Proof.} Homomorphisms inject into homomorphisms of their
integral first homology lattices, so \(\Lambda_i\) is a finitely
generated free abelian group. Pullback on invariant one-forms depends
linearly on the homomorphism. Writing a polarization as a positive sum
of \(i\eta\wedge\bar\eta\) shows that \(E_i\) is a positive semidefinite
real quadratic form. There is no nonzero integral vector in its kernel:
if all pulled-back one-forms vanish on \(Y\), the map \(u\alpha\) has
zero derivative and is constant; because \(\alpha(Y)\) generates \(B\)
and both maps preserve the chosen origins, \(u=0\). The intersection
expression for \(E_i\) is integral, so polarization gives a rational
Gram matrix in every integral basis. A nonzero real kernel of a rational
matrix contains a nonzero rational vector, and clearing its denominator
would give a nonzero integral kernel vector. Thus the real kernel is
zero and \(q_i\) is positive definite. Its positive values on nonzero
lattice vectors are at least one.

Let \(\lambda_1\le\cdots\le\lambda_{r_i}\) be the successive minima of
the Euclidean unit ball for this metric. Minkowski's second theorem
gives {[} v\_\{r\_i\}\prod\_j\lambda\_j\le 2\^{}\{r\_i\}C\_i. {]} Since
every \(\lambda_j\ge\sqrt{\mu_i}\), one obtains
\(\lambda_{r_i}\le R_i\). The minima are attained, giving the required
independent maps. They need not be an integral basis, and none is
asserted: a rational spanning set is enough. If \(P\) vanished on every
tuple from these spanning sets, multilinearity would make it vanish
identically, a contradiction. A map with image dimension smaller than
\(d\) pulls every smooth real \(2d\)-form back to zero, so the nonzero
wedge period gives the image-dimension assertion.

For the degree bound, write \(\beta=F^*c_1(L)\), a nef class. Its energy
is \(E=\beta H^{d-1}=\sum_iE_i(u_{ij_i})\le\sum_iR_i^2\). The
mixed-intersection inequality {[}
\beta\^{}d\le\frac{(\beta H^{d-1})^d}{(H^d)^{d-1}} \tag{B10.2} {]}
follows directly from surface Hodge index: for rational ample
\(\beta,H\), take smooth complete-intersection surfaces of suitable very
ample multiples of \(j-1\) copies of \(\beta\) and \(d-j-1\) copies of
\(H\), for \(1\le j\le d-1\). Hodge index on each surface gives
\(a_j^2\ge a_{j-1}a_{j+1}\) for \(a_j=\beta^jH^{d-j}\). These
inequalities imply \(a_d\le a_1^d/a_0^{d-1}\). Approximating nef
\(\beta\) by \(\beta+\varepsilon H\) proves the required nef case by
continuity. For \(d=1\) it is equality, and for \(d=2\) it is B5.
Projection formula gives \(\beta^d=\delta_F\deg_L Z\), proving (B10.1).
Finally nefness of \(Bc_1(L)-c_1(L')\) gives \(E_{L'}\le B E_L\) after
pullback and intersection with \(H^{d-1}\), so the same argument proves
the last assertion. QED.

\textbf{Exact imported lattice source.} Martin Henk, \emph{Successive
minima and lattice points}, arXiv:math/0204158v1, Theorem 1.3, printed
p.2, and its proof in §3, printed pp.5--7,
\url{https://arxiv.org/pdf/math/0204158}. The statement and complete
proof were read; only the displayed upper-product inequality is imported
here.

\textbf{What determinant is required.} On the real Hom space, let
\(w_1,\ldots,w_r\) be a basis of normalized analytic vectors, with real
Gram matrix \(G\). If the columns of \(B\) express an integral basis of
\(\Lambda\) in these coordinates, then exactly {[}
C\_\Lambda\textsuperscript{2=\det(B}\{\mathsf t\}GB)=\textbar{}\det B\textbar\^{}2\det G.
\tag{B10.3} {]} Replacing the analytic basis by an orthonormal one
replaces \(B\) by \(G^{1/2}B\) and preserves this covolume. Rallis can
contribute to \(G\) after the CM-period and measure comparisons are
fixed; it does not identify \(B\). If \(\Gamma=\sum\mathbb Zw_j\) is
commensurable with \(\Lambda\), then {[} \frac{C_\Lambda}{C_\Gamma}
=\frac{[\Gamma:\Gamma\cap\Lambda]}{[\Lambda:\Gamma\cap\Lambda]}. {]}
Thus a proved integral analytic sublattice \(\Gamma\subset\Lambda\) with
bounded Gram determinant already gives \(C_\Lambda\le C_\Gamma\); no
saturation-index upper bound is then needed. If instead
\(\Lambda\subset\Gamma\), its index contributes in the numerator and
must be bounded. A uniform integer \(M\) with \(M\Gamma\subset\Lambda\)
gives the sufficient bound \(C_\Lambda\le M^r\sqrt{\det G}\). This
identifies the exact direction of the integral comparison that a
theta-index theorem would have to supply.

\textbf{Variation and normalization.} The ranks \(r_i\) may grow when
genuinely new cohomology appears at deeper source levels. Neither a
rank-independent constant nor a uniform covolume is supplied by finite
generation. For the fixed sublattice of old maps pulled through a
degree-\(e\) cover of the source, pull back \(H\) as well. Then
\(E_i\mapsto eE_i\), \(C_i\mapsto e^{r_i/2}C_i\), and a known \(\mu_i\)
becomes \(e\mu_i\). Consequently \(R_i\mapsto\sqrt e R_i\),
\(V\mapsto eV\), and \(\delta_F\mapsto e\delta_F\); the reduced-image
bound (B10.1) is unchanged. The actual generic degree in that bound must
be retained. Replacing it by the trivial lower bound one can lose an
essential cancellation. This proposition gives a sufficient determinant
criterion for producing a small witness from prescribed source lattices.
It is not a necessary condition for intrinsic bounded witness degree.

\subsection{B11. A fixed-order CM Hecke example separates the Hom
regulator from witness
degree}\label{b11.-a-fixed-order-cm-hecke-example-separates-the-hom-regulator-from-witness-degree}

Let {[} E=\mathbb C/(\mathbb Z+i\mathbb Z),\qquad
E\_m=\mathbb C/(\mathbb Z+mi\mathbb Z),\quad m\ge1, {]} with their
degree-one polarizations. Then {[}
\operatorname{Hom}(E,E\_m)=m\mathbb Z{[}i{]},\qquad
\deg(m(a+bi))=m(a\textsuperscript{2+b}2). \tag{B11.1} {]} In particular
the least positive degree grows with \(m\), despite fixed rank and fixed
CM field.

\textbf{Proof.} A homomorphism is multiplication by a complex number
\(\gamma\) with \(\gamma\mathbb Z[i]\subset\mathbb Z+mi\mathbb Z\).
Write \(\gamma=x+mi y\) with \(x,y\in\mathbb Z\). The further
requirement \(i\gamma\in\mathbb Z+mi\mathbb Z\) forces \(m\mid x\). Thus
\(\gamma=m(a+bi)\), and the converse is immediate. The degree is the
index of \(\gamma\mathbb Z[i]\) in the target lattice, namely
\(|\gamma|^2/m=m(a^2+b^2)\). These elliptic curves are algebraic and
have CM; multiplication by \(m\) gives a degree-\(m\) isogeny, so they
lie in one polarized Hecke orbit. QED.

Here is a version with a fixed integral quadratic action and actual Weil
witnesses. Put \(A_m=E_m^4\) with product principal polarization
\(L_m\), and let \(\mathbb Z[i]\) act by sending \(i\) to {[}
J=\operatorname{diag}\left(

\begin{pmatrix}0&-1\\1&0\end{pmatrix}

,

\begin{pmatrix}0&-1\\1&0\end{pmatrix}

\right). {]} This is an integral action for every \(m\), has Rosati
involution \(J^\dagger=-J\), and has multiplicities \((2,2)\) on
\(H^{1,0}\), so \(A_m\) is of Weil type for the same order. The diagonal
isogeny \(E^4\to E_m^4\) given by multiplication by \(m\) commutes with
this action and pulls the target polarization back to \(m\) times the
source polarization. Thus the example stays in the corresponding
polarized quadratic-action Hecke orbit.

Let \(Y=E^2\), \(H\) its product degree-one polarization, and consider
{[} f\_m(x,y)=(mx,0,my,0),\qquad
Z\_m=E\_m\times\{0\}\times E\_m\times\{0\}\subset A\_m. {]} The image
has \(\deg_{L_m}Z_m=2\). Its Weil projection is nonzero: the extreme
\(i\)-eigenline contains the nonzero \((2,2)\)-form {[}
(dz\_1+i,dz\_2)\wedge(dz\_3+i,dz\_4) \wedge(d\bar z\_1+i,d\bar z\_2)
\wedge(d\bar z\_3+i,d\bar z\_4), {]} whose restriction to \(Z_m\) is
\(dz_1\wedge dz_3\wedge d\bar z_1\wedge d\bar z_3\) and has nonzero
integral. Pairing with this Weil form detects the nonzero projection, by
the Poincaré self-adjoint projector in B6.

The full fixed-source Hom lattice is {[}
\operatorname{Hom}(E\^{}2,A\_m)=M\_\{4\times2\}(m\mathbb Z{[}i{]}), {]}
of rank \(16\). In the integral basis furnished by the real and
imaginary matrix units, its energy metric is \(m I_{16}\): the
polarization form on \(E_m\) is \(\frac{i}{2m}dz\wedge d\bar z\), so
multiplication by \(m\gamma\) contributes \(m|\gamma|^2\) to the energy.
Consequently its least nonzero energy is \(m\), its Gram determinant is
\(m^{16}\), and its metric covolume is \(m^8\). Orthonormalizing the
analytic coordinates transfers the factor \(\sqrt m\) to the integral
coefficients and changes none of these intrinsic metric quantities.

For the selected two-dimensional map, however, {[}
f\_m\^{}*c\_1(L\_m)=mH,\quad
E\_\{f\_m\}=2m,\quad V=H\^{}2=2,\quad \delta\emph{\{f\_m\}=m\^{}2, {]}
and therefore {[} \frac{E_{f_m}^2}{\delta_{f_m}V}=2=\deg}\{L\_m\}Z\_m.
{]} All assertions follow by the lattice-index calculation, product
intersection, and the displayed invariant form. This supplies a precise
counterexample to any claim that finite generation or an orthonormal
analytic basis uniformly bounds integral map energy on a CM Hecke orbit.
At the same time it supplies bounded nonzero Weil witnesses. It
therefore does \textbf{not} refute intrinsic bounded division; it shows
that a raw Hom-regulator bound is sufficient but unnecessary, and that
generic-degree cancellation is mathematically essential.

\subsection{B12. The exact geometric quantity is an infimum of
normalized
energies}\label{b12.-the-exact-geometric-quantity-is-an-infimum-of-normalized-energies}

For a polarized Weil-type abelian \(2n\)-fold \((A,L)\), let
\(D_W(A,L)\) be the least degree of an effective integral \(n\)-cycle
with nonzero Weil projection, or \(+\infty\) if none exists. Then {[}
D\_W(A,L)=\inf\_\{(Y,f,H)\} \frac{\bigl(f^*c_1(L)\cdot H^{n-1}\bigr)^n}
\{\delta\_f (H\textsuperscript{n)}\{n-1\}\}, \tag{B12.1} {]} where \(Y\)
runs over smooth projective \(n\)-folds, \(f:Y\to A\) over morphisms
generically finite onto an image \(Z\) with nonzero Weil projection,
\(\delta_f\) is the generic mapping degree, and \(H\) runs over integral
ample classes on \(Y\). The infimum of the empty set is understood to be
\(+\infty\).

\textbf{Proof.} The mixed-intersection inequality B10.2 and projection
formula show that every expression on the right is at least
\(\deg_L Z\ge D_W(A,L)\). If a witness exists, at least one irreducible
component has nonzero Weil projection and no larger degree. Positive
degrees are integers, so a least-degree irreducible witness \(Z\)
exists. Take a projective resolution \(f:Y\to Z\subset A\), with
\(\delta_f=1\), and put \(\beta=f^*c_1(L)\). This is integral, big and
nef, and \(\beta^n=\deg_L Z=D_W(A,L)\). Fix an integral ample \(H_0\).
The classes \(H_j=j\beta+H_0\) are integral ample. As \(j\to\infty\),
{[}
\beta H\_j\textsuperscript{\{n-1\}=j}\{n-1\}D\_W+O(j\^{}\{n-2\}),\qquad
H\_j\textsuperscript{n=j}nD\_W+O(j\^{}\{n-1\}). {]} For \(n=1\) the
expression is already equality; for \(n\ge2\) the quotient in B12.1
tends to \(D_W\). This proves the reverse inequality and hence equality.
QED.

This is an exact variational characterization, not an effective
construction. It explains what a successful automorphic estimate would
have to control: normalized energies for some witnesses, allowing their
actual mapping degrees and source volumes. Requiring every theta
coefficient or the entire Hom regulator to remain bounded is stronger
than necessary. B10 provides one workable sufficient criterion with
explicit constants. B11 shows why it is not an equivalence on a
prescribed source. The infimum in B12 allows resolutions and
polarizations of the unknown witnesses, so it does not close the open
existence or uniformity problem by itself.

\subsection{Computation attached to
B6}\label{computation-attached-to-b6}

\texttt{weil\_projector\_certificate.py} constructs the projector for
\(K=\mathbb Q(\omega)\), \(\omega^2+\omega+1=0\), \(a=3+\omega\) (trace
5, norm 7), for \(n=2,3,4\), and checks exact polynomial idempotence
modulo \(f_W f_R\). It also checks the norm-product resultant formula.
This is a computation checking those examples, not a formal proof of B6
for all \(n\). For \(n=2\), for example, {[}
f\_W=X\textsuperscript{2-23X+2401,\quad f\_R=X}3-126X\^{}2+6174X-117649,
{]} {[} R=25,729,836,300. {]} Reducing rational polynomial coefficients
gives a smaller possible denominator \(476\,478\,450\) in this example.
Minimality here means the least common denominator of this unique
degree-\(\leq4\) interpolation polynomial, not a claim of a minimal
denominator among all algebraic projector correspondences.

\subsection{Review of approach paths}\label{review-of-approach-paths}

\begin{enumerate}
\def\labelenumi{\arabic{enumi}.}
\tightlist
\item
  Tried direct GRH bound at the asserted edge. Refined by checking the
  actual Rallis dimensions: genuine torus-to-three-dimensional lift is
  at \(3/2\), and B1 is stronger and unconditional for the global
  factor.
\item
  Tried transfer from theta norm to degree. Direct identity is not
  justified; B5 supplies a valid inequality conditional on controlling
  all actual polarization coordinate pullbacks.
\item
  Tried conductor control along principal levels. B3 gives exact
  presentation growth and shows conductor is not level; no inference to
  intrinsic \(c(s)\) is made.
\item
  Tried use of bounded arbitrary projector norms. B9 refutes it.
  Replaced it by the fixed polynomial projector B6 in the generic
  \(K\)-action.
\item
  Tried treating the resulting rank-two coefficient index as a separate
  unbounded gate. B7 and B8 remove it once an effective degree bound is
  given, with explicit quadratic bounds. The residual gate is uniformly
  bounded effective witnesses with nonzero Weil component, or
  deformation of a full class that supplies them.
\end{enumerate}

\subsection{Suggested exact Grothendieck placement, to be independently
source-checked by
root}\label{suggested-exact-grothendieck-placement-to-be-independently-source-checked-by-root}

B6 concerns algebraic correspondences and primitive projectors in
homological/numerical motives; B7 uses the Hodge standard positivity and
the primitive Lefschetz decomposition; B8 concerns integral lattices
underlying those motives together with proper cycle parameter spaces.
The attached sketch points to cote 16 scans 44--50 for the
standard-conjecture correspondence framework. No original archive scans
are supplied here; this report does not claim to have read or corrected
their transcription. Root must cite exact EGA/SGA/FGA and archive
locators only after actual verification.

\subsection{Additional checked classical sources for B5 and
B7}\label{additional-checked-classical-sources-for-b5-and-b7}

\begin{itemize}
\tightlist
\item
  Eduardo Cattani, \emph{Introduction to Kähler manifolds}, ICTP June
  2010 preliminary notes, §5.3, Proposition 5.15 and Theorem 5.17,
  printed p.38: the primitive Hodge-star formula and Hodge--Riemann
  positivity. Example 5.19 and equation (5.23), printed pp.39--40,
  specialize to the surface Hodge inequality. The displayed statements
  and the adjoining positivity proof were read. Source:
  \url{https://people.maths.ox.ac.uk/joyce/Nairobi2019/CattaniKahlerManifolds.pdf}.
  These are classical imported results; this report does not claim an
  independent proof of the full Hodge--Riemann theorem.
\item
  Jean-Pierre Demailly, \emph{Complex Analytic and Differential
  Geometry}, version 21 June 2012, Chapter III, §1.D, equation (1.24)
  and Wirtinger's inequality (1.25), printed pp.136--137: analytic-cycle
  volume and its linear-algebra proof. Source:
  \url{https://people.math.harvard.edu/~demarco/Math274/Demailly_ComplexAnalyticDiffGeom.pdf}.
  Demailly uses \(\omega_D=i\sum h_{jk}dz_j\wedge d\bar z_k\) and the
  real metric \(\operatorname{Re}h\), giving \(\omega_D^p/(2^p p!)\).
  For the same real metric our convention is \(\omega=\omega_D/2\), so
  this is precisely \(\omega^p/p!\) as used in B7. The formula and proof
  were read and the factor of two checked.
\end{itemize}

\section{Semiregularity and the replacement-object problem}

% Candidate statements and complete proofs for the DEG_L paper.
% Checked by /root/semireg_t7, GPT-6 Astra Pro; inherited effort not exposed.
% This is an input fragment, not a standalone article.
% Source locators are written explicitly so no guessed bibliography keys are used.

\subsection{The actual surface and its deformation problem}

\begin{proposition}[S-01: Allowed source covers do not preserve an embedding premise]
\label{prop:semireg-cover-domain}
Let $f:Y\to A$ be a surface morphism obtained from the four rational
Hodge maps in 032, Proposition~3.1. If $\pi:Y'\to Y$ is a connected finite
\'etale cover of degree $d>1$, the pulled-back maps still satisfy its
nonzero-period criterion, but $f\pi$ is not a closed immersion.
Moreover $(f\pi)_*[Y']=d f_*[Y]$.
\end{proposition}
\begin{proof}
Pullback preserves rational Hodge maps and multiplies the integral of
a top-degree class by $d$. Every geometric point of $Y$ has $d$ distinct
preimages with the same image under $f\pi$, so this morphism is not
injective on geometric points. The last assertion is the proper
pushforward formula for a finite map. The deeper principal covers
used here are explicitly permitted in 032, Corollary~6.7 and \S7.
\end{proof}

\begin{proposition}[S-02: The exact conormal gate]
\label{prop:semireg-dual-gate}
For a smooth closed immersion $i:Y\hookrightarrow A$ of a surface in
an abelian $g$-fold, $g\geq3$, write $N=N_{Y/A}$ and
$W=H^0(A,\Omega_A^1)$. Bloch semiregularity is equivalent to
surjectivity of
\[
 H^1(A,\mathcal O_A)\otimes\bigwedge^3W
 \longrightarrow H^1(Y,N^\vee\otimes K_Y).
\]
The map is restriction followed by the natural morphism
$\bigwedge^3 i^*\Omega_A^1\to N^\vee\otimes K_Y$.
\end{proposition}
\begin{proof}
Bloch's map is
$H^1(Y,N)\to H^{g-1}(A,\Omega_A^{g-3})$.
Serre duality identifies its dual source with $H^1(A,\Omega_A^3)$
and its dual target with $H^1(Y,N^\vee\otimes K_Y)$.
The conormal exact sequence induces the displayed exterior-power
map, since the quotient $\Omega_Y^1$ has rank two.
This is the defining dual map in Buchweitz--Flenner \S8.1;
Proposition~8.2 identifies it with Bloch's map.
Finally $\Omega_A^3=\mathcal O_A\otimes\bigwedge^3W$,
and finite-dimensional duality gives the equivalence.
\end{proof}

\begin{proposition}[S-03: A nonzero period does not force semiregularity]
\label{prop:semireg-period-counterexample}
There is a smooth surface in a CM abelian threefold with a nonzero
product of two holomorphic and two antiholomorphic one-forms for
which the Bloch semiregularity kernel has dimension at least eight.
\end{proposition}
\begin{proof}
Let $E$ be a CM elliptic curve, $B=E^2$, and
$L=\mathcal O_E(3[0])\boxtimes\mathcal O_E(3[0])$.
A general $C\in|L|$ is smooth and connected. Since $L^2=18$ and
$K_B=\mathcal O_B$, adjunction gives $g(C)=10$ and
$N_{C/B}=K_C$. Put $A=B\times E$ and $Y=C\times E$.
Then $N_{Y/A}=K_C\boxtimes\mathcal O_E$, so K\"unneth gives
$h^1(Y,N)=1+10=11$. The divisor semiregularity target
$H^2(A,\mathcal O_A)$ has dimension three.
The curve $C$ dominates the first elliptic factor: otherwise its
ample divisor class would be a fibre class. Hence a nonzero one-form
from that factor restricts to a nonzero form $\alpha$ on $C$.
For a nonzero one-form $\beta$ on the last $E$, the integral
$\int_Y(\alpha\wedge\beta)\wedge\overline{\alpha\wedge\beta}$
is nonzero, by the product formula and positivity of the two
one-dimensional integrals.
\end{proof}

\begin{proposition}[S-04: Rigid generating sources]
\label{prop:semireg-rigid-source}
Let $Y$ be a smooth projective surface with $H^1(Y,T_Y)=0$ and
$i:Y\hookrightarrow A$ a smooth closed immersion. Let $G\subset A$
be the abelian subvariety generated by a translate of $i(Y)$.
Then the ambient obstruction map satisfies
\[
\ker(\operatorname{ob}_i)=
\operatorname{Lie}(A)\otimes
\ker\bigl(H^1(A,\mathcal O_A)\to H^1(Y,\mathcal O_Y)\bigr).
\]
In particular it is injective when $G=A$. If in addition $i$ is
Bloch semiregular, this is also the infinitesimal Hodge-locus kernel
of the entire class $i_*[Y]$.
\end{proposition}
\begin{proof}
Restriction on $H^1(T_A)$ is $\operatorname{id}_{\operatorname{Lie}(A)}
\otimes i^*$, since $T_A$ is trivial. The normal sequence and
$H^1(T_Y)=0$ make $H^1(i^*T_A)\to H^1(N)$ injective.
This proves the kernel formula. The map $\operatorname{Alb}(Y)\to A$
has image $G$, so $i^*$ on $H^1(\mathcal O)$ is injective precisely
when $G=A$. Buchweitz--Flenner, Corollary~4.12 and Proposition~8.2,
identify
\[
 \sigma_i(\operatorname{ob}_i(\kappa))
 =\kappa\mathbin{\lrcorner}i_*[Y].
\]
If $\sigma_i$ is injective, the kernels of the two sides agree.
\end{proof}

The actual 032 source is a compact smooth complex two-ball quotient.
Calabi--Vesentini, \emph{Sur les vari\'et\'es complexes compactes
localement sym\'etriques}, Th\'eor\`eme~3, p.~315, gives
$H^1(Y,T_Y)=0$. The nonzero four-factor period does not establish
$G=A$. For example, separate injectivity on the two eigenspaces
of a quadratic action is insufficient: in
$A=E^n\times E^n$ with action $(a,\bar a)$, pullback to the diagonal
$E^n$ is injective on each eigenspace, but their images coincide.

\begin{proposition}[S-05: A quantitative obstruction for the actual morphism]
\label{prop:semireg-map-dimension}
Let $Y$ be a compact smooth complex two-ball quotient and
$f:Y\to A$ a morphism with two-dimensional image, where $\dim A=g$.
Let $\tau_f$ be the Buchweitz--Flenner morphism semiregularity map.
Then
\[
 \dim\ker\tau_f\ \geq\ gq(Y)-\binom{2g}{g-2}.
\]
\end{proposition}
\begin{proof}
In characteristic zero $df:T_Y\to f^*T_A=\mathcal O_Y^{\oplus g}$
is generically injective. Its kernel is a torsion-free sheaf of
generic rank zero, hence vanishes. With $N_f=\operatorname{coker}(df)$,
the derived dual of $L_{Y/A}$ is the two-term complex
$[T_Y\to f^*T_A]$ in degrees $0,1$, so
$T^2_{Y/A}(\mathcal O_Y)=H^1(Y,N_f)$.
Calabi--Vesentini and the cokernel sequence give
$H^1(\mathcal O_Y)^{\oplus g}\hookrightarrow H^1(N_f)$.
The target specified in Buchweitz--Flenner, Proposition~7.21,
Lemma~7.22 and the definition before Theorem~7.23, is
\[
 \bigoplus_{r=0}^{g-2}H^{r+2}(A,\Omega_A^r).
\]
Its dimension is
$\sum_{r=0}^{g-2}\binom gr\binom g{r+2}=\binom{2g}{g-2}$
by Vandermonde. Rank--nullity proves the assertion.
\end{proof}

\begin{corollary}[S-06: Failure at sufficiently deep allowed levels]
\label{cor:semireg-level-failure}
Fix a nonzero 032 surface construction $f:Y\to A$. There are
connected finite \'etale covers $\pi_j:Y_j\to Y$, preserving its
four pulled-back Hodge maps and nonzero period, for which
$\tau_{f\pi_j}$ is not injective. Such covers may be chosen among
a cofinal system of principal congruence levels.
\end{corollary}
\begin{proof}
The nonzero period implies $H^1(Y,\mathbf C)\ne0$.
Venkataramana, \emph{Virtual Betti numbers of compact locally
symmetric spaces}, Theorem~1, proves that finite-dimensional
direct-limit cohomology over finite covers equals the cohomology
of the compact dual. Here the compact dual is $\mathbf P^2$,
whose first cohomology is zero. Pullback under finite covers is
injective over $\mathbf Q$ by transfer, so the direct limit is
nonzero and therefore infinite-dimensional. Hence $q(Y_j)$ is
unbounded, and the preceding proposition applies with fixed $A$.
The congruence version is explicitly identified in the final Remark
on p.~3 of that paper. Its proof uses the same rational-conjugation
action, since rational conjugation and finite intersection preserve
congruence subgroups; principal subgroups are cofinal.
Persistence of the period follows from the first proposition.
\end{proof}

This is a statement about the deformation problem of a morphism.
The factorization
$\tau_f=\sigma_{Rf_*\mathcal O_Y}\circ\psi_f$
does not imply that a kernel of $\tau_f$ survives in the perfect
complex obstruction space: $\psi_f$ itself can have a kernel.
Nor does the corollary settle Bloch semiregularity of a particular
smooth reduced image.

\begin{proposition}[S-07: Descent of the cover obstruction space]
\label{prop:semireg-cover-invariants}
For a finite \'etale Galois cover $\pi:Y'\to Y$ with deck group $H$,
$N_{f\pi}=\pi^*N_f$ and
$H^1(Y',N_{f\pi})^H=H^1(Y,N_f)$.
\end{proposition}
\begin{proof}
Flat pullback identifies the tangent cokernel sequences.
Finite pushforward and the projection formula identify
$\pi_*\pi^*N_f=N_f\otimes\pi_*\mathcal O_{Y'}$.
In characteristic zero averaging over $H$ splits off precisely
the invariant summand $N_f$. Cohomology commutes with this direct
summand, giving the assertion.
\end{proof}

\subsection{Derived replacements and their remaining gate}

\begin{proposition}[S-09: An integral Chern-component projector]
\label{prop:semireg-k-projector}
Let $A$ be an abelian $g$-fold and $\xi\in K_0(A)$ with
$\operatorname{ch}_j(\xi)=0$ for $j<p$. For $m\geq2$ put
\[
 M=[m]^*,\qquad
 D=\prod_{j=p+1}^g(m^{2j}-m^{2p}),\qquad
 \eta=\prod_{j=p+1}^g(m^{2j}\operatorname{id}-M)\xi.
\]
Then $\eta\in K_0(A)$ and
$\operatorname{ch}(\eta)=D\operatorname{ch}_p(\xi)$ in rational
Betti cohomology.
\end{proposition}
\begin{proof}
The operator $M$ multiplies $H^{2j}(A,\mathbf Q)$ by $m^{2j}$.
Thus every higher Chern component is annihilated by its corresponding
factor, lower components were already zero, and the component of
index $p$ is multiplied by $D$. The polynomial has integer
coefficients. Integral K-classes are represented by perfect complexes
on the smooth projective variety, using finite locally free
resolutions and shifts to implement negative coefficients.
\end{proof}

Starting from an integral codimension-$p$ cycle of class $N_0t$,
one may use $\xi=\sum a_i[\mathcal O_{Z_i}]$; its leading Chern
character is the fundamental cycle. Alternatively, resolve each
$Z_i$ and use the proper pushforward of its structure sheaf;
Grothendieck--Riemann--Roch gives the same leading component.
Adding line-bundle classes gives the useful character
$Nt+\sum_a b_a e^{a\theta}$, which remains Hodge on the polarized
Route-2 family. This constructs a K-class, not a semiregular object.

\begin{proposition}[S-08: Multiplicity creates trace kernels]
\label{prop:semireg-multiplicity}
If $E$ is perfect and $\operatorname{Ext}^2(E,E)\ne0$, then
$E^{\oplus m}$ is not fully semiregular for any $m\geq2$.
\end{proposition}
\begin{proof}
Its degree-two self-Ext group is
$\operatorname{Ext}^2(E,E)\otimes\operatorname{End}(\mathbf C^m)$.
The Atiyah class is $\operatorname{at}(E)\otimes\operatorname{id}$;
therefore every semiregularity component is the corresponding
component for $E$ tensored with the matrix trace.
A nonzero degree-two Ext class tensored with a nonzero traceless
matrix lies in the kernel.
\end{proof}

\begin{proposition}[S-10: Full Fourier--Mukai transport]
\label{prop:semireg-full-fm}
For the fixed enhanced Fourier--Mukai equivalence in the supplied
032 companion \S5.2, the full trace maps satisfy
\[
 \sigma_{\Phi Q}\circ\Phi_{\rm Ext}
   =\Phi^H_{HH}\circ\sigma_Q
 \quad\hbox{on }\operatorname{Ext}^2_A(Q,Q).
\]
Thus full semiregularity is invariant under this equivalence.
\end{proposition}
\begin{proof}
Use the companion's enhanced adjunctions and notation.
Put
$b_Q(v)=\operatorname{Tr}_{Q,H_A[2]}(A_Q^H[2]\circ v)
\in HH_{-2}(A)$.
Equations (5.20), (5.24) identify its HKR image with $\sigma_Q(v)$.
For $\delta\in T^{-2}(A)$, trace-compatible adjunction gives
\[
 \langle b_Q(v),\delta\rangle
   =\operatorname{Tr}_{A,Q}(\delta_Q[2]\circ v).
\]
Indeed, if $\delta':H_A\to\omega_A[g-2]$ is the adjoint coefficient
map, move $\delta'[2]$ through the partial trace by (5.24);
evaluation then gives the right side. This is exactly (5.37)
with the displayed even shift.
The convolution proof of (5.33), with the same inverse adjunction
data and shifted by two, gives
$(\Psi^{-2}\delta)_{\Phi Q}=\theta_Q[-2]\Phi(\delta_Q)$.
By the trace preservation in (5.31),
\[
 \langle b_{\Phi Q}(\Phi_{\rm Ext}v),\Psi^{-2}\delta\rangle
  =\operatorname{Tr}_{A,Q}(\delta_Q[2]\circ v)
  =\langle b_Q(v),\delta\rangle.
\]
Definition (5.34) and nondegeneracy of the Serre pairing imply
$b_{\Phi Q}(\Phi_{\rm Ext}v)=\Phi_*^{HH}b_Q(v)$.
Apply (5.35) and the fixed HKR maps. The two vertical maps
are isomorphisms, proving the assertion about kernels.
This proves a refinement; Proposition~5.2 itself only asserted
transport on the obstruction-evaluation image.
\end{proof}

\begin{proposition}[S-11: An ample complete intersection need not be semiregular]
\label{prop:semireg-ci-counterexample}
Let $A$ be an abelian fourfold, $L$ very ample, and
$Y=D_1\cap D_2$ a smooth transverse intersection of two members
of $|L|$. Then $h^1(N_{Y/A})=24$, whereas
$h^3(A,\Omega_A^1)=16$. Thus its Bloch map has a kernel of
dimension at least eight.
\end{proposition}
\begin{proof}
The twisted Koszul sequence is
$0\to L^{-1}\to\mathcal O_A^{\oplus2}\to\mathcal I_Y\otimes L\to0$.
Kodaira vanishing and $K_A=\mathcal O_A$ give $H^j(A,L)=0$ for
$j>0$; Serre duality gives $H^j(A,L^{-1})=0$ for $j<4$.
Thus
$H^2(A,\mathcal I_Y\otimes L)=H^2(A,\mathcal O_A)^{\oplus2}$.
The restriction sequence gives
$H^1(Y,L|_Y)=H^2(A,\mathcal I_Y\otimes L)$, of dimension twelve.
Transversality identifies $N_{Y/A}=(L|_Y)^{\oplus2}$.
The Bloch target has dimension
$\binom43\binom41=16$, so rank--nullity finishes the proof.
\end{proof}

\begin{proposition}[S-12: A pure Weil character obstruction for coherent sheaves]
\label{prop:semireg-pure-weil-no-go}
Let $A_0$ be a projective member of Voisin's four-dimensional
marked $\mathbf Q(i)$-Weil torus family, and let $t\ne0$ be a
rational Weil class. No fully Buchweitz--Flenner-semiregular
coherent sheaf $F$ on $A_0$ satisfies
\[
 \operatorname{ch}_1(F)=\operatorname{ch}_3(F)=0,\qquad
 \operatorname{ch}_2(F)=t.
\]
The rank and top-degree character are unrestricted.
\end{proposition}
\begin{proof}
The unpolarized Weil torus parameter space is smooth near $A_0$;
Voisin, \emph{A counterexample to the Hodge conjecture extended
to K\"ahler varieties}, \S3, equations (3.1)--(3.8), describes it.
Lemma~3.1 shows that the rational Weil space remains of Hodge
type $(2,2)$. Hence the components with indices $I=\{1,2,3\}$
in the assertion remain horizontal Hodge classes. In dimension
four these are precisely the components tested by full
semiregularity.
Buchweitz--Flenner Theorem~5.1 and its proof on p.~179 would
therefore extend $F$ as a coherent flat sheaf on a neighborhood.
Voisin Proposition~3.2 and Proposition~2.1 show that very general
fibres in that neighborhood have $\operatorname{NS}=0$ and
$c_2(G)=0$ for every coherent analytic sheaf $G$.
On such a fibre $c_1(F)=c_2(F)=0$, whence
$\operatorname{ch}_2(F)=0$, contrary to the nonzero horizontal
class $t$.
\end{proof}

This last proposition is a coherent-sheaf statement for the
specified fourfold family. Extension to arbitrary perfect complexes
would require an analytic effectivity theorem for their formal
deformations; algebraic/formal Pridham alone is not that theorem.
Allowing a nonzero polarization term in the Chern character
avoids this particular contradiction.

\section{Bounded division in families}

Run: DEG\_L, 2026-11-21. Author: \texttt{/root}, GPT-6 Astra Pro; effort
setting not exposed. These are new formulations and proofs for this run.
They have not yet received the independent whole-proof check requested
below. Equality of a cycle with a rational cohomology class always means
equality under the cycle-class map, not equality in a Chow group whose
right-hand side is undefined.

\subsection{DIV1. A proper parameter space for a bounded signed
representative}\label{div1.-a-proper-parameter-space-for-a-bounded-signed-representative}

\textbf{Setup.} Let \texttt{f:X\ -\textgreater{}\ S} be a smooth
projective morphism over a smooth irreducible complex algebraic variety
of finite type. Its fibres have dimension \texttt{g}. Fix a relatively
very ample line bundle \texttt{H}, an integer
\texttt{0\ \textless{}=\ p\ \textless{}=\ g}, and a global flat section
\texttt{t} of \texttt{R\^{}\{2p\}f\_*\ Q(p)}. Put \texttt{r=g-p}. For a
complex point \texttt{s}, let \texttt{c\_t(s)} be the least integer
\texttt{C\ \textgreater{}=\ 1} for which there are an integer
\texttt{1\ \textless{}=\ N\ \textless{}=\ C} and effective integral
\texttt{r}-cycles \texttt{Z+}, \texttt{Z-} on \texttt{X\_s}, each of
\texttt{H}-degree at most \texttt{C}, such that

\texttt{N\ t\_s\ =\ cl(Z+)\ -\ cl(Z-)}.

Set \texttt{c\_t(s)=infinity} if no such representative exists. Zero
cycles in the sense of the zero element of the cycle group are allowed
among the effective cycles, including when \texttt{r\textgreater{}0}.

\textbf{Statement.} For each integer \texttt{C}, the subset
\texttt{S\_C=\{s\ :\ c\_t(s)\textless{}=C\}} is Zariski closed.
Consequently:

\begin{enumerate}
\def\labelenumi{\arabic{enumi}.}
\tightlist
\item
  A bound on any Zariski-dense subset gives the same bound on all of
  \texttt{S}.
\item
  If \texttt{t\_s} is algebraic at every complex point, then
  \texttt{S\_C=S} for some \texttt{C}.
\item
  The same conclusion holds if \texttt{t\_s} is algebraic outside a
  countable union of proper Zariski-closed subsets.
\end{enumerate}

Thus algebraicity on the entire family is equivalent to bounded signed
integral complexity on a dense subset. This conclusion supplies
existence of a bound, not an algorithm for computing it.

\textbf{Proof.} We spell out the finiteness which the informal reference
to a bounded-degree Chow scheme suppresses. Grothendieck,
\emph{Techniques de construction et théorèmes d'existence en géométrie
algébrique IV: Les schémas de Hilbert}, Séminaire Bourbaki, exposé 221
(1960/61), Lemme 2.4, printed p.255, states that the reduced,
pure-dimensional closed subschemes of fibres of a projective family with
bounded degree form a bounded family. Applying it in dimension
\texttt{r} and degree at most \texttt{C}, and stratifying a finite-type
parameter space by flatness, gives a finite list of Hilbert polynomials
\texttt{P\_1,...,P\_k} containing those of all reduced irreducible
subvarieties which can occur as components of these cycles. Each
corresponding relative Hilbert scheme is projective over \texttt{S}, by
the same exposé, Théorème 3.2, printed p.260.

Every nonzero effective cycle of degree at most \texttt{C} is a sum
\texttt{sum\_i\ m\_i{[}V\_i{]}}, where \texttt{V\_i} is reduced and
irreducible of dimension \texttt{r}, \texttt{m\_i} is a positive
integer, there are at most \texttt{C} summands, and
\texttt{m\_i\textless{}=C}. Indeed every component has positive integral
degree. There are only finitely many lists of the polynomials and
multiplicities with total degree at most \texttt{C}. For each list take
the corresponding fibre product over \texttt{S} of Hilbert schemes.
Include \texttt{S} itself to parameterize the zero cycle. This gives
finitely many projective parameter schemes whose universal
top-dimensional fundamental cycles include every desired effective
cycle. Additional points with nonreduced schemes or lower-dimensional
embedded components cause no problem: the top-dimensional fundamental
cycle is still effective and its degree is the leading coefficient of
the fixed Hilbert polynomial, multiplied by \texttt{r!}.

Take pairs of these parameter schemes for the two signs. On each such
finite-type projective parameter scheme \texttt{T\ -\textgreater{}\ S},
the class of the difference of the two universal fundamental cycles is a
flat section of the pulled-back cohomology local system. The following
argument verifies this also when \texttt{T} is singular or nonreduced.
For each universal subscheme \texttt{Z\ -\textgreater{}\ T}, its
structure sheaf is coherent and \texttt{T}-flat, and the ambient
morphism \texttt{X\_T\ -\textgreater{}\ T} is smooth. Thus \texttt{O\_Z}
is an absolutely perfect complex on \texttt{X\_T}: locally the ambient
ring is flat over the base with regular fibres of bounded dimension, and
a base-flat coherent module has finite Tor amplitude over that ambient
ring. Precisely this assertion is Stacks Project, Lemma 15.79.5, tag
\texttt{068X}, with base Tor amplitude \texttt{{[}0,0{]}}. It is
absolute perfection on \texttt{X\_T}, rather than only relative
perfection over \texttt{T}, that is used here.

One may work over affine opens of \texttt{S}; then \texttt{X\_T} is
quasi-projective, so the perfect complex admits a bounded locally free
resolution and defines a topological Chern character on
\texttt{X\_T(C)}. Derived restriction to any fibre is the ordinary
\texttt{O\_\{Z\_t\}} because of \texttt{T}-flatness. Its component
\texttt{ch\_p(O\_\{Z\_t\})} equals the Betti class of the
top-dimensional fundamental cycle \texttt{{[}Z\_t{]}\_r}, including its
multiplicities. An exact algebraic locator is Fulton, \emph{Intersection
Theory}, second edition, Theorem 18.3(3),(5), printed p.353, together
with Example 18.3.11, printed p.363: the latter identifies the top term
of the Riemann--Roch class with the associated cycle of any coherent
sheaf; on a smooth ambient fibre the Riemann--Roch class is
\texttt{ch(O\_\{Z\_t\})\ td(TX\_t)\ cap\ {[}X\_t{]}}, and the Todd
factor has constant term one. Lower-dimensional support components do
not contribute in codimension \texttt{p}.

Therefore all the fibre classes in question are restrictions of one
topological class \texttt{ch\_p(O\_Z)} on \texttt{X\_T(C)}. The smooth
proper family \texttt{X\ -\textgreater{}\ S} is locally a differentiable
fibre bundle; its pullback to \texttt{T(C)} is topologically locally
trivial. Restricting a global cohomology class in such a family gives a
flat section, since the fibre inclusions over any path in a trivializing
neighborhood are homotopic. Complex algebraic spaces are locally path
connected, and their finitely many connected components agree in the
analytic and algebraic topologies. The argument also treats nilpotents
because passing to the reduction does not alter the underlying analytic
topological space. Taking sums and differences gives the asserted flat
section for the universal signed cycle.

For any fixed \texttt{N} between 1 and \texttt{C}, equality of this
section with \texttt{N\ t} therefore selects a union of connected
components of \texttt{T}. There are finitely many such components, each
open and closed. The selected union is projective over \texttt{S}. Its
image is Zariski closed. A finite union over the parameter schemes and
\texttt{N} is precisely \texttt{S\_C}, proving the first assertion. A
closed subset containing a Zariski-dense subset is all of \texttt{S},
which proves (1).

Algebraicity of a rational class at a point is equivalent to finite
\texttt{c\_t(s)}: clear all coefficients of an algebraic rational cycle
and split the resulting integral cycle by sign. Hence algebraicity at
all points says \texttt{S(C)=union\_\{C\textgreater{}=1\}\ S\_C(C)}. If
every \texttt{S\_C} were proper, choose a small closed analytic ball
with nonempty interior in \texttt{S(C)}. Each \texttt{S\_C} meets that
ball in a closed nowhere-dense set, since a proper algebraic subset of
an irreducible smooth complex variety contains no analytic open subset.
The Baire category theorem rules out a countable covering of the ball by
these sets. Thus some \texttt{S\_C=S}. The same argument, adding the
given countably many exceptional proper closed subsets to the purported
covering, proves (3). QED.

\textbf{Scope and Grothendieck locator.} This is a statement about
relative algebraic cycles and their Betti realization. It uses FGA 221,
Lemme 2.4 and Théorème 3.2 for boundedness and properness. The analogous
cycle classes in any compatible Weil realization are transported by the
same algebraic correspondences, but a new assertion of integral Tate
surjectivity is not part of the statement. It proves the precise form
needed from sketch Route 2, statement 2.1; it does not use André--Oort
or any asserted density of a particular Hecke orbit.

\subsection{DIV2. Bounded witnesses on a dense set suffice in a Weil
family}\label{div2.-bounded-witnesses-on-a-dense-set-suffice-in-a-weil-family}

\textbf{Statement.} Let \texttt{S} be a smooth irreducible complex
algebraic variety of finite type, as in DIV1. Let
\texttt{A\ -\textgreater{}\ S} be a smooth projective family of
polarized abelian varieties of dimension \texttt{2n} and Weil type for a
fixed imaginary quadratic order, with globally defined order action,
fixed polarization type and a fixed rational flat Weil class \texttt{t}.
A finite base change may be made to define the action and the class.
Assume there is a constant \texttt{D\_0} and a Zariski-dense subset
\texttt{T} such that for every \texttt{s\ in\ T} there is an effective
integral \texttt{n}-cycle \texttt{Z\_s} of degree at most \texttt{D\_0}
whose Weil component is nonzero. Then \texttt{c\_t} is bounded on all of
\texttt{S}. An explicit bound on \texttt{T}, and therefore on
\texttt{S}, is the ceiling of the maximum of the two quadratic
expressions in B8 evaluated at \texttt{D=D\_0}, together with 1.

\textbf{Proof.} Apply B6--B8 with the fixed generic quadratic action.
Their projector denominator, signed-degree multiplier, polarization
volume and Weil-lattice discriminant are constant on the family; so is
the norm of a fixed flat multiple of \texttt{t}. B8 produces, for each
witness, an integral multiple of \texttt{t\_s} and signed effective
representatives with the displayed uniform quadratic bounds. This gives
\texttt{T\ subset\ S\_C} for one integer \texttt{C}. DIV1 gives
\texttt{S\_C=S}. If the polarization is only relatively ample, replace
it by a fixed relatively very ample power; degrees change by the fixed
factor equal to that power raised to \texttt{n}, so boundedness is
equivalent. QED.

The conclusion removes the need to control denominators of a projector
defined using each point's larger CM field. It leaves a geometric
requirement: uniformly bounded effective witnesses with a nonzero
component. It does not prove that requirement from the existence of
rational cycles on each CM fibre.

\subsection{DIV3. What a successful semiregular cycle construction would
prove}\label{div3.-what-a-successful-semiregular-cycle-construction-would-prove}

\textbf{Statement.} In the setting of DIV1, with
\texttt{1\ \textless{}=\ p\ \textless{}=\ g}, suppose at one point
\texttt{s\_0} there are smooth regularly embedded subvarieties
\texttt{Y\_1,...,Y\_k} of codimension \texttt{p} and fixed integers
\texttt{a\_j,N}, \texttt{N\textgreater{}=1}, with
\texttt{N\ t\_\{s\_0\}=sum\_j\ a\_j\ cl(Y\_j)}. Suppose each embedding
is Bloch semiregular and the flat continuation of \textbf{each full
class} \texttt{cl(Y\_j)} is of Hodge type on an analytic neighborhood of
\texttt{s\_0}. Then \texttt{c\_t} is bounded on all of \texttt{S}.

\textbf{Proof.} Apply Buchweitz--Flenner, \emph{A semiregularity map for
modules and applications to deformations}, Compositio Math.137 (2003),
Theorem 5.2, p.175, with the singleton set of indices \texttt{I=\{p\}}.
The hypotheses are a deformation of a compact complex algebraic central
manifold over a smooth complex germ, a horizontal Hodge section
\texttt{alpha\_p}, and an \texttt{I}-semiregular subspace with
\texttt{alpha\_p(s\_0)=ch\_p(O\_Y)}. Our smooth projective family over
the analytic germ of \texttt{S} meets the ambient hypotheses. For a
regular codimension-\texttt{p} immersion,
\texttt{T\^{}2\_\{Y/X\}(O\_Y)=H\^{}1(Y,N\_\{Y/X\})},
\texttt{ch\_p(O\_Y)={[}Y{]}}, and Proposition 8.2, p.200, identifies the
component \texttt{tau\_\{p-1\}} with Bloch's map. Thus Bloch
semiregularity is exactly the required \texttt{I}-semiregularity, and
preservation of this full fundamental class is sufficient; no other
component of \texttt{ch(O\_Y)} is required.

The proof of Theorem 5.2, p.180, invokes the same lifting argument as
the proof of Theorem 5.1, p.179, using the subspace obstruction maps of
§§4.11 and 4.13 in place of the sheaf maps. That argument constructs
formal lifts over the infinitesimal neighborhoods of the smooth base
germ, uses a versal deformation, and applies Artin approximation to
obtain an actual local flat deformation. For subspaces this is a local
section of the relative Douady space, equivalently of the
analytification of the relative Hilbert scheme in the present projective
setting. Apply it separately to every \texttt{Y\_j} and intersect the
finitely many parameter neighborhoods. The local extensions have the
original Hilbert polynomials.

The degrees of the fundamental cycles of these extensions are constant,
because their Hilbert polynomials are fixed. Their cohomology classes
are the flat continuations of the original classes, so the displayed
signed equality with \texttt{N\ t} persists. Splitting the fixed
integers \texttt{a\_j} by sign gives effective representatives of
degrees at most \texttt{sum\_j\ \textbar{}a\_j\textbar{}\ deg\_H(Y\_j)},
uniformly on a nonempty analytic open set. It is Zariski dense in
\texttt{S}. DIV1 therefore extends the bound to every fibre.
Equivalently, the finite product of the relevant projective Hilbert
schemes has a component dominating that analytic open set; its proper
image is all of \texttt{S}. QED.

\textbf{Necessary distinctions.} Preserving only
\texttt{sum\_j\ a\_j\ cl(Y\_j)} does not meet the stated hypotheses.
Semiregularity of a perfect complex is likewise a statement about its
Atiyah-trace map and the Chern components detected by that map; it is
not supplied by its \texttt{K\_0} class. A projective bounded parameter
space or an independent local-to-cycle argument is needed when a sheaf
or complex replaces the embeddings. S-01--S-07 show why the original
rigid 032 ball quotient does not directly satisfy these conditions on a
varying generated abelian family.

\subsection{DIV4. An explicit upper bound for the natural Hecke
representative}\label{div4.-an-explicit-upper-bound-for-the-natural-hecke-representative}

\textbf{Statement.} Let \texttt{psi:A\_x\ -\textgreater{}\ A\_0} be an
isogeny of polarized abelian \texttt{g}-folds, with equal polarization
type and \texttt{psi\^{}*L\_0=m\ L\_x} at the level of first Chern
classes, where \texttt{m} is a positive integer. Suppose
\texttt{psi\^{}*\ t\_0=(a/b)t\_x}, with coprime integers \texttt{a!=0},
\texttt{b\textgreater{}0}. If
\texttt{N\_0t\_0=cl(Z\_0\^{}+)-cl(Z\_0\^{}-)} and both degrees are at
most \texttt{D\_0}, then

\texttt{c\_t(x)\ \textless{}=\ max(N\_0\ \textbar{}a\textbar{},\ b\ m\^{}p\ D\_0)}.

\textbf{Proof.} Isogenies in characteristic zero are finite flat, so
flat pullback of each effective integral cycle is effective integral.
Pulling back the given equality and multiplying by \texttt{b} gives
\texttt{N\_0\ a\ t\_x=b(psi\^{}*cl(Z\_0\^{}+)-psi\^{}*cl(Z\_0\^{}-))}.
Reverse the signs if \texttt{a\textless{}0}. B4 computes the degree of
each pullback as \texttt{m\^{}p} times the old degree. This gives
exactly the bound asserted. QED.

This upper bound is for a specified representative. It is not a lower
bound for the minimum \texttt{c\_t(x)}. When \texttt{a/b=m\^{}p},
normalizing the rational pullback keeps its numerical degree fixed while
its displayed denominator grows. T6.4--T6.5 cancel auxiliary source
covers and target multiplication for specially normalized correspondence
constructions; they do not cancel an arbitrary Hecke transport between
different target fibres.

\subsection{DIV5. What the known fourfold theorem already
supplies}\label{div5.-what-the-known-fourfold-theorem-already-supplies}

\textbf{Statement.} On any fixed finite-type family of complex abelian
fourfolds and any global rational flat Hodge class, the signed integral
complexity is uniformly bounded, provided one admits the supplied
Markman theorem of the rational Hodge conjecture for all abelian
fourfolds. The same conclusion applies to any higher-dimensional family
for which the supplied theorem proves algebraicity of the particular
flat class on every fibre.

\textbf{Proof.} The admitted theorem supplies algebraicity pointwise. If
the base is smooth and irreducible, DIV1(2) supplies a uniform bound.
For a general finite-type base, choose a finite stratification of its
reduction into smooth irreducible locally closed complex subvarieties,
apply DIV1(2) to the restricted family and flat class on each stratum,
and take the maximum of the finitely many resulting bounds. Fibres over
the original base and its reduction are identical at complex points.
QED.

This bound is a posteriori and ineffective in this argument. It does not
calculate a theta degree, identify a Rallis factor, or extend the known
fourfold theorem to the open Weil families in sketch T8. The exact
supplied Markman locators are Theorem 1.5.1 (\texttt{thm-algebraicity}),
proving algebraicity of the Hodge--Weil classes on polarized abelian
sixfolds of Weil type for an arbitrary imaginary quadratic field and
discriminant \texttt{-1}, and Corollary 1.6.1, asserting the Hodge
conjecture for all abelian fourfolds. These numbers were checked from
the source's subsection-based shared theorem counter after removing
comments and hidden draft blocks; the sixfold assertion must not be
broadened to arbitrary discriminant or arbitrary Hodge classes.

\clearpage
\part{Open questions and unproved implications}
\section{The remaining main questions}
No statement in this part is asserted as a theorem without an explicit
conditional antecedent. The exact raw degree is not a function of source
level alone. A complete normalization specifies the primitive integral
maps, their polarization, theta inputs, adelic measures and local factors.
The finite spectral formula does not automatically reduce every rank-two
standard $L$-function to Hecke characters.

The uniform geometric question is the existence, on dense fibres of each
open Weil family, of effective middle-dimensional cycles of bounded degree
whose Weil component is nonzero. B6--B8 and DIV1--DIV2 remove the generic
projector and bounded-division extraction from that question. They do not
prove the bounded witness assertion. The integral Hom covolume, its rank,
and the generic mapping degree must be retained when converting theta
norms into a bound for the reduced image.

For a specified smooth reduced 032 image the Bloch question remains the
surjectivity in S-02. The universal assertion for all allowed deeper-level
map presentations is refuted by S-05--S-06. A deformation of a projected
class requires a different semiregular object, or another argument; the
full unprojected rigid surface class need not stay Hodge. A projected
$K_0$ class is not a proof that a representative object is semiregular.

Every approach tried, its outcome, and the next precise gate are recorded
in the component reports, QUESTIONS.tsv, and LEDGER.md. The scheduled
three-hour progress review has not yet occurred in this checkpoint.
\section{Consequences and exact framework locators}

This file is a live research deliverable. The session remains in
progress; open questions below have not been promoted to theorems. All
findings use run ID \texttt{DEG\_L\_20261121}; their check events are
collected in \texttt{CHECK\_EVENTS.tsv}.

\subsection{The degree and L-value
question}\label{the-degree-and-l-value-question}

T6.1 identifies the effective surface degree with a positive sum of
norms of all pulled-back polarization-coordinate wedges. T6.2 separates
it from the pure four-factor codimension-two output, whose degree for
the product polarization is zero. T6.3 refutes a raw degree determined
by source level alone, because the construction permits multiplication
of its denominator-clearing maps at the same level. The meaningful
normalized Chow cycle is unchanged under auxiliary covers and these map
multiplications once its divisor kernels are chosen as biextensions
(T6.4--T6.5).

T6.7 removes a proposed analytic obstruction: algebraic norm and trace
clear the nonzero extracted period without requiring a lower bound for
its individual complex absolute value. T6.8 bounds the effective degrees
created by fixed divisor kernels. These statements retain the
coefficient field, integral generator lattice and kernel constants
explicitly.

T6.9 uses the supplied human proof and the checked theta spanning
literature to cover every pulled-back holomorphic coordinate. The
Hodge-index estimate B5 therefore reduces the effective surface degree
to a finite character-theta Gram form. The global character factor
occurs at 3/2 and has uniform elementary Euler-product bounds for fixed
field degree (B1). Uniformity of integral maps, local factors and
normalization is a separate issue. The exact wedge-norm computation
through a rank-two source is being pursued; no character-only
factorization has been assumed.

\subsection{Route 2 and the Hodge
conjecture}\label{route-2-and-the-hodge-conjecture}

B6 constructs the Weil projector using only a fixed generic imaginary
quadratic action, with an explicit resultant denominator. B7 bounds the
projected component of an effective cycle by its degree in a positive,
flat Weil metric. B8 turns any nonzero effective Weil witness into a
representative of every fixed Weil class, with both denominator and
signed effective degree bounded quadratically in the witness degree.
DIV1 supplies the proper parameter-space argument; DIV2 carries a
uniform witness bound on a dense set to every fibre.

This is a strict reduction of the denominator and projector problem. It
is not a proof of uniformly bounded witnesses in the open Weil families.
No new general Hodge-conjecture case is asserted from the presently
established estimates. The supplied Markman theorem yields a posteriori
bounded complexity for abelian fourfold families by DIV5. That existence
argument gives no numerical constant or explicit theta representative.

\subsection{Semiregularity and
deformation}\label{semiregularity-and-deformation}

The 032 paper supplies a morphism with two-dimensional image, not a
smooth embedded surface. S-01 shows that its permitted covers do not
even preserve injectivity of that morphism. S-02 is the exact
conormal-cohomology test for an actual smooth embedding. S-03 and S-11
refute general shortcuts from a nonzero period or from an ample complete
intersection to semiregularity.

For the rigid ball-quotient source, S-04 forces the generated quotient
to remain constant in any embedded semiregular deformation of its entire
class. S-05--S-06 refute universal semiregularity of all deeper-level
morphism presentations even for the full generalized map: the
obstruction space grows beyond its fixed Hodge target. S-07 removes that
cover growth by equivariant descent but returns to the base-level
obstruction.

S-09 produces a perfect-complex K-class whose surviving Chern character
is the desired projected class plus a chosen polarization polynomial. It
removes extraneous Hodge conditions at the K-class level. It does not
construct a semiregular object. S-10 explains why an equivalence alone
does not remove an existing full trace kernel. S-12 shows that
preserving a pure Weil character without polarization data can force a
contradiction with Voisin's nonprojective torus family for coherent
sheaves. The remaining geometric problem is a specific object with an
injective obstruction trace and compatible full Hodge data. DIV3 states
precisely how a successful embedded construction would imply bounded
division.

\subsection{GRH hypotheses}\label{grh-hypotheses}

The task's Alpöge Dirichlet GRH theorem is admitted, with the supplied
Lean declarations as its exact formal statement. Consequently the cited
conditional Dirichlet edge estimates may be used without a further GRH
assumption in this run. The full Lean dependency archive was not
supplied or compiled.

General Hecke edge estimates are conditional on the Masullo claim
identified by the prompt as §1.2. The supplied index does not contain
the complete claim or proof and does not identify a numbered Hecke-GRH
theorem. Its numbered trace formulas have not been substituted for such
a statement. No Hecke GRH hypothesis is needed for the absolutely
convergent 3/2 factor. A nonabelian rank-two standard L-function is
outside the stated Dirichlet/Hecke scope unless a factorization is
established.

\subsection{Tate and standard
conjectures}\label{tate-and-standard-conjectures}

The supplied 032 and Milne consequences remain the background theorems
requested by the task: CM Hodge, Tate for abelian varieties over finite
fields in the cited formulation, and the Hodge standard conjecture for
abelian varieties. The new degree estimates do not enlarge their scope.
The correspondences in B6 exist algebraically and define projectors in
homological realizations and numerical motives; their cohomological
calculation alone does not claim idempotence in the Chow ring. A
positive cohomological norm does not by itself construct an effective
cycle or control an integral saturation index.

\subsection{Grothendieck locators and the previous
runs}\label{grothendieck-locators-and-the-previous-runs}

\begin{itemize}
\tightlist
\item
  \textbf{FGA, Bourbaki exposé 221, Lemme 2.4 p.255 and Théorème 3.2
  p.260:} bounded reduced subvarieties and projective
  fixed-Hilbert-polynomial parameter spaces give DIV1--DIV3 their
  precise algebraic-geometric mechanism.
\item
  \textbf{SGA 7 I, Exposé VII, (2.4.2), Example 2.9.5 and Remark 2.9.6:}
  biextension additivity gives the exact cancellation T6.5. Original
  scan verification is being completed.
\item
  \textbf{SGA 7 I, Exposé VI:} the deformation framework belongs to
  obstruction theory, with the precise semiregularity maps supplied by
  Bloch and Buchweitz--Flenner, not by a positivity assertion in the
  archive.
\item
  \textbf{Cote 16, the Coates letter, PDF page 45 / archive footer 44,
  and PDF page 46 / footer 45:} the displayed equivalence involving
  algebraic Lefschetz inverse and cycle-rank symmetry is consistent with
  retaining the algebraic projector as a separate datum. The current
  results do not infer missing cycles from positivity alone.
\item
  \textbf{Cote 44, PDF page 2 / archive footer 1, Th.4; PDF page 4 /
  footer 3, Th.5:} polarization and the intersection/height form remain
  distinct from the unnormalized degree of the 032 surface. No new claim
  about the arithmetic height formula is made in this run.
\end{itemize}

The sketch uses PDF-page numbering, including the cover, for several of
these scan references. We record both PDF and archive-footer numbers;
this is a locator convention, not a mathematical transcription
correction.

The rows directly treated are \texttt{HODGE\_BRIDGE\_SKETCH\_20261121}
T6, T7, Route 2 §2.1, and the generic-projector/witness part of §2.4. T8
is affected only by the reduction DIV2; its higher-dimensional cases
remain open here. Other runs named SA06, SA07, H7, H8 and H9 were not
supplied beyond quotations in the sketch; no unseen row is declared
proved or changed.


\end{document}
